\documentclass{article}
\usepackage{amsmath, amsthm, amssymb, appendix, bm, graphicx, hyperref, mathrsfs, tikz-cd, enumerate,natbib}
%\usepackage[nocap]{ctex}
%\usepackage{ctex}
\usepackage[english]{babel}

\title{\textbf{Deformations of Calabi-Yau varieties}}
\author{hwn}
\date{\today}
\linespread{1.5}
\theoremstyle{definition}
\newtheorem{definition}{Definition}[subsection]
\newtheorem{example}[definition]{Example}
\newtheorem{remark}[definition]{Remark}
\theoremstyle{plain}
\newtheorem{theorem}[definition]{Theorem}
\newtheorem{lemma}[definition]{Lemma}
\newtheorem{corollary}[definition]{Corollary}
\newtheorem{proposition}[definition]{Proposition}
\renewcommand{\abstractname}{\Large\textbf{Abstract}}
\usepackage{geometry}
\geometry{left=2.54cm,right=2.54cm,top=2.54cm,bottom=2.54cm}
\setcounter{secnumdepth}{5}

\begin{document}

\pagenumbering{arabic}
\maketitle
\begin{abstract}
    This article gives a systematic introduction to the deformation theory of Calabi–Yau varieties through two complementary approaches: complex geometry and algebra.
    In the complex‑geometric part, we start from the Maurer–Cartan equation, establish the Kodaira–Spencer map,
    and discuss in detail the role of the Tian–Todorov lemma in solving the equation, deriving the unobstructedness theorem for smooth Calabi‑Yau manifolds.
    In the algebraic part, we construct the Lichtenbaum–Schlessinger \(T^i\) functors and the cotangent complex,
    present the obstruction theory for deformations of schemes, prove Schlessinger’s criterion and pro‑representability of the deformation functor,
    and then introduce the Ran–Kawamata \(T^1\) lifting principle together with its application to smooth Calabi‑Yau varieties.
    At the intersection of the two theories, the final part is devoted to threefold Calabi–Yau varieties with ordinary double points.
    We describe Friedman’s smoothability criterion and the subsequent generalization due independently to Kawamata, Ran and Tian:
    as long as a Kähler resolution exists, deformations are unobstructed, so Friedman’s first condition can be dropped.
\end{abstract}
\tableofcontents
\newpage

\section{Deformation theory with approaches of complex geometry}

\subsection{The Maurer-Cartan equation}

A complex manifold is by definition a differentiable manifold $ M $ endowed with an atlas whose transition functions are holomorphic. As is known to all, a complex structure defined by such an atlas can also be encoded by an integrable almost complex structure $ I $, which is just an endomorphism of the tangent bundle satisfying a certain non-linear integrability condition. In this section, we will work with the latter description of complex structures. We will explain how to view the various complex structures $ I $ on $ M $ or, more precisely, the (small) deformations of a given one.

We start out with the set\[
    \mathcal{A}_{ac}(M):=\{I\mid I^2=-\mathrm{id}\}\subset \mathrm{End}(TM)
\]of all almost complex structures on $ M $. First consider an arbitrary $ I\in \mathcal{A}_{ac}(M) $. Recall that an almost complex structure $ I $ is uniquely determined by a decomposition of the tangent bundle $ T_{\mathbb{C} }M=T^{1,0}\oplus T^{0,1} $ with $ I=i\cdot \mathrm{id} $ on $ T^{1,0} $ and $ I=-i\cdot\mathrm{id} $ on $ T^{0,1} $. In fact, giving $ T^{0,1}\subset T_{\mathbb{C} }M $ is enough, for $ T^{1,0}=\overline{T^{0,1}} $.

If $ I(t) $ is a continuous family of almost complex structures with $ I(0)=I $, one has a continuous family of such decompositions $ T_{\mathbb{C} }M=T^{1,0}_t\oplus T^{0,1}_t $ or, equivalently, of subspaces $ T^{0,1}_t\subset T_{\mathbb{C} }M $.

Thus, for small $ t $ the deformation $ I(t) $ of $ I $ can be encoded by a map\[
    \phi(t):T^{0,1}\rightarrow T^{1,0}\ \text{with}\ v+\phi(t)(v)\in T^{0,1}_t.
\]We write $ T^{1,0} $ and $ T^{0,1} $ for the subbundles defined by $ I $.
% Explicitly, one has\[
%     \phi(t)=-\mathrm{pr}_{T^{1,0}_t}\circ j,
% \]where $ j:T^{0,1}\subset T_{\mathbb{C} } $ and $ \mathrm{pr}_{T^{1,0}_t}:T_{\mathbb{C} }\rightarrow T^{1,0}_t $ are the natural inclusion respectively projection.

Conversely, if $ \phi(t) $ is given, then one defines for small $ t $\[
    T^{0,1}_t:=(\mathrm{id}+\phi(t))(T^{0,1}).
\]The condition $ t $ be `small' has to be imposed in order to ensure that with this definition $ T^{0,1}_t $ has the right dimension. More precisely, $ t $ small means that $ T^{0,1}_t\subset T_{\mathbb{C} }\rightarrow T^{0,1} $ is an isomorphism, which is certainly an open condition, at least on a compact manifold. To avoid problems caused by non-compactness we will henceforth assume that $ M $ is compact.

Let us now consider the power series expansion\[
    \phi(t)=\phi_0+\phi_1 t+\phi_2 t^2+\cdots
\]of a given $ \phi $. Then $ \phi_0=0 $ and the higher order coefficients $ \phi_{i>0} $ will be regarded as sections of $ \bigwedge^{0,1}\otimes T^{1,0} $.

Next, we shall study the integrability of the deformed almost complex structures $ I(t) $. From now on we will assume that $ I $ is integrable and we will denote the complex manifold $ (M,I) $ by $ X $. In particular, the higher order coefficients $ \phi_i $ correspond to elements in $ \mathcal{A}^{0,1}(\mathcal{T}_X) $.

The integrability condition for $ I(t) $ can be expressed as\[
    [T^{0,1}_t,T^{0,1}_t]\subset T^{0,1}_t.
\]

Let us rephrase this in terms of $ \phi(t) $. First note that there exists the $ \bar{\partial} $-operator on the holomorphic tangent bundle $ \mathcal{T}_X $ of $ X=(M,I) $, which can in particular be applied to $ \phi(t)\in \mathcal{A}^{0,1}(\mathcal{T}_X) $. Secondly, one defines a `Lie bracket'\[
    [\ ,\ ]:\mathcal{A}^{0,p}(\mathcal{T}_X)\times \mathcal{A}^{0,q}(\mathcal{T}_X)\longrightarrow \mathcal{A}^{0,p+q}(\mathcal{T}_X)
\]by taking the Lie bracket in $ \mathcal{T}_X $ and the exterior product in the form part $ \mathcal{A}^{0,*} $. Spelled out in local coordinates this gives for $ \alpha=\sum d\bar{z}_I\otimes v_I $ and $ \beta=\sum d\bar{z}_J\otimes w_J $ with $ v_I=\sum \alpha_{Ij}\frac{\partial}{\partial z_j} $ and $ w_J=\sum \beta_{Jl}\frac{\partial}{\partial z_l} $ the following formula\[
    [\alpha,\beta]=\sum (d\bar{z}_I \wedge d\bar{z}_J)[v_I , w_J]=\sum_{I,j,J,l}(d\bar{z}_I \wedge d\bar{z}_J)\big[\alpha_{Ij}\frac{\partial}{\partial z_j},\beta_{Jl}\frac{\partial}{\partial z_{l}}\big].
\]This is indeed well-defined, as any holomorphic coordinate change would change the tangent vector fields $ v_I,w_I $ by antiholomorphic factors $ \bar{f} $ and $ \frac{\partial\bar{f}}{\partial z_m}=0 $.

\begin{remark}\label{cvb}
    Using the local description, one easily checks that $ \bar{\partial}[\alpha,\beta]=[\bar{\partial}\alpha,\beta]\pm[\alpha,\bar{\partial}\beta] $. In particular, the Lie bracket of two $ \bar{\partial} $-closed sections in $ \mathcal{A}^{0,*}(\mathcal{T}_X) $ is again closed and the bracket of a closed one and an exact one is exact. Thus, the bracket induces a map\[
        H^p(X,\mathcal{T}_X)\times H^q(X,\mathcal{T}_X)\longrightarrow H^{p+q}(X,\mathcal{T}_X).
    \]
\end{remark}

\begin{lemma}
    The integrability equation $ [T^{0,1}_t,T^{0,1}_t]\subset T^{0,1}_t $ is equivalent to the \textbf{Maurer-Cartan equation}\begin{equation}\label{wer}
        \bar{\partial}\phi(t)+[\phi(t),\phi(t)]=0.
    \end{equation}
\end{lemma}

\begin{proof}
    The assertion can be proved locally. We will omit the parameter $ t $ in the notation and shall write\[
        \phi=\sum_{i,j}\phi_{ij}d\bar{z}_i\otimes \frac{\partial}{\partial z_j}.
    \]If the integrability condition holds true, then we have in particular\[
        \bigg[\frac{\partial}{\partial \bar{z}_i}+\phi\bigg(\frac{\partial}{\partial \bar{z}_i}\bigg),\frac{\partial}{\partial \bar{z}_k}+\phi\bigg(\frac{\partial}{\partial\bar{z}_k}\bigg)\bigg]\in T^{0,1}_t.
    \]Since $ [\frac{\partial}{\partial \bar{z}_i},\frac{\partial}{\partial \bar{z}_k}]=0 $, this is equivalent to\[
        \biggl(\sum_l \bigg[\frac{\partial}{\partial \bar{z}_i},\phi_{kl}\frac{\partial}{\partial z_l}\bigg]+\sum_j\bigg[\phi_{ij}\frac{\partial}{\partial z_j},\frac{\partial}{\partial\bar{z}_k}\bigg]+\sum_{j,l}\bigg[\phi_{ij}\frac{\partial}{\partial z_j},\phi_{kl}\frac{\partial}{\partial z_l}\bigg] \biggr)\in T^{0,1}_t.
    \]Now,\[
        \bigg[\frac{\partial}{\partial\bar{z}_i},\phi_{kl}\frac{\partial}{\partial z_l}\bigg]=\frac{\partial \phi_{kl}}{\partial\bar{z}_i}\frac{\partial}{\partial z_l}\ \text{and}\ \bigg[\phi_{ij}\frac{\partial}{\partial z_j},\frac{\partial}{\partial\bar{z}_k}\bigg]=-\frac{\partial\phi_{ij}}{\partial \bar{z}_k}\frac{\partial}{\partial z_j}
    \]and, hence,\begin{align*}
        &\sum_l\bigg[\frac{\partial}{\partial\bar{z}_i},\phi_{kl}\frac{\partial}{\partial z_l}\bigg]+\sum_j\bigg[\phi_{ij}\frac{\partial}{\partial z_j},\frac{\partial}{\partial \bar{z}_k}\bigg]\\
        =&\sum_{j}\bigg(\frac{\partial\phi_{kj}}{\partial \bar{z}_i}-\frac{\partial\phi_{ij}}{\partial\bar{z}_k}\bigg)\frac{\partial}{\partial z_j}\\
        =&(\bar{\partial}\phi)\bigg(\frac{\partial}{\partial\bar{z}_i},\frac{\partial}{\partial\bar{z}_k}\bigg).
    \end{align*}On the other hand,\[
        \sum_{j,l}\bigg[\phi_{ij}\frac{\partial}{\partial z_j},\phi\frac{\partial}{\partial z_l}\bigg]=[\phi,\phi]=\bigg(\frac{\partial}{\partial \bar{z}_i},\frac{\partial}{\partial\bar{z}_k}\bigg).
    \]Therefore, the integrability condition implies that $ \bar{\partial}\phi+[\phi,\phi]\in\mathcal{A}^{0,2}(\mathcal{T}^{1,0}X) $ satisfies\[
        (\bar{\partial}\phi+[\phi,\phi])\bigg(\frac{\partial}{\partial\bar{z}_i},\frac{\partial}{\partial\bar{z}_k}\bigg)\in T^{0,1}_t.
    \]Hence, $ \bar{\partial}+[\phi,\phi]\in\mathcal{A}^{0,2}(T^{1,0}X\cap T^{0,1}_t) $. But for small $ t $ one has $ T^{1,0}X\cap T^{0,1}_t =0 $ and, therefore, $ \bar{\partial}\phi+[\phi,\phi]=0 $.

    Coversely, if the Maurer-Cartan equation holds true, then integrability condition is satisfied for a local frame of $ T^{0,1}_t $. This immediately shows that it is true for all sections of $ T^{0,1}_t $.
\end{proof}

Let us now consider the power series expansion of the Maurer-Cartan equation, i.e.\ we replace $ \phi $ in (\ref{wer}) by $ \phi=\sum_{i=1}^{\infty}\phi_i t^i $:\[
    \bar{\partial}\bigg(\sum_{i=1}^{\infty}\phi_i t^i\bigg)+\sum_{i,j=1}^{\infty}[\phi_i,\phi_j]t^{i+j}=0.
\]This yields a recursive system of equations:\begin{align*}
    0&=\bar{\partial}\phi_1\\
    0&=\bar{\partial}\phi_2+[\phi_1,\phi_1]\\
    \vdots& \\
    0&=\bar{\partial}\phi_k +\sum_{0<i<k}[\phi_i,\phi_{k-i}].
\end{align*}

We shall first study the first-order equation $ \bar{\partial}\phi_1=0 $ in detail. The Maurer-Cartan equation is, in particular, saying that the first-order deformation of the complex structure $ I $ is described by a $ \bar{\partial} $-closed $ (0,1) $-form $ \phi_1 $ with values in the holomorphic tangent bundle $ \mathcal{T}_X $. Thus, it defines an element $ [\phi_1]\in H^1(X,\mathcal{T}_X) $.

\begin{definition}
    The \textbf{Kodaira-Spencer class} of a one-parameter deformation $ I_t $ of the complex structure $ I $ is the induced cohomology class $ [\phi_1]\in H^1(X,\mathcal{T}_X) $.
\end{definition}

In order to identify isomorphic (first order) deformations, the (infinitesimal) action of the diffeomorphism group $ \mathrm{Diff}(M) $ has to be taken into account. Let $ F_t $ be a one-parameter family of diffeomorphisms of the manifold $ M $. With respect to local coordinate functions $ x_i $ this can be expanded into a power series $ \tilde{x_i}=x_i+tF_i (x)+t^2\dots $. The first order term $ \sum_i F_i \frac{\partial}{\partial x_i} $ defines a global vector field\[
    \frac{dF_t}{dt}|_{t=0}\in \mathcal{A}^0(TM).
\]

We can rewrite this identification on the complex manifold $ X=(M,I) $ withh respect to holomorphic coordinates $ z_i $ as $ \tilde{z_i}=z_i+tf_i(z,\bar{z})+t^2\dots $. For the antiholomorphic coordinates $ \bar{z_i} $ we then have $ \bar{\tilde{z_i}}=\bar{z_i}+t\bar{f_i}(z,\bar{z})+t^2\dots $. Thus, $ \frac{dF_t}{dt}|_{t=0}\in\mathcal{A}^0(TX) $ in local coordinates (up to the factor 2) is given as\[
    \sum f_i\frac{\partial}{\partial z_i}+\sum \bar{f}_i\frac{\partial}{\partial\bar{z}_i}.
\]

Any diffeomorphism can be used to construct new complex structures out of the given one $ I $ by pushing-forward $ I $ via the differential of the diffeomorphism. More explicitly, for the diffeomorphism $ F_t $ the new complex structure is given by $ dF_t\circ I\circ {(dF_{t})}^{-1} $. The antiholomorphic tangent bundle of this complex structure is thus\[
    T^{0,1}_t=dF_t(T^{0,1}X).
\]

In local coordinates the first order deformation of $ T^{0,1}X $ is described as the image of $ T^{0,1}X $ under\[
    \mathrm{id}+\sum_{i,j}\frac{\partial f_i}{\partial z_j}\otimes\frac{\partial}{\partial z_j}+\frac{\partial f_i}{\partial\bar{z_j}}d\bar{z_j}\otimes\frac{\partial}{\partial z_i}+\frac{\partial \bar{f_i}}{\partial z_j}dz_j\otimes\frac{\partial}{\partial \bar{z_i}}+\frac{\partial\bar{f_i}}{\partial\bar{z_j}}d\bar{z_j}\otimes\frac{\partial}{\partial \bar{z_i}}.
\]In this sum, only the third and fifth terms are non-trivial on $ T^{0,1}X $. Moreover, the fifth term is an endomorphism of $ T^{0,1}X $.

Thus, the first order deformation of $ T^{0,1}X $ induced by a one-parameter family of diffeomorphisms $ F_t $ of $ M $ is described by the homomorphism\[
    T^{0,1}X\longrightarrow T^{1,0}X,\ \frac{\partial}{\partial\bar{z_j}}\longmapsto \sum_i \frac{\partial f_i}{\partial \bar{z_j}}\frac{\partial}{\partial z_i}.
\]This yields

\begin{lemma}
    The first-order deformation of the complex structure $ I $ on the manifold $ M $ (defining the complex manifold $ X=(M,I) $) induced by a one-parameter family $ F_t $ of diffeomorphisms of $ M $ is determined by the map\[
        \phi_1=\bar{\partial}\bigg({\big(\frac{dF_t}{dt}|_{t=0}\big)}^{1,0}\bigg):T^{0,1}X\longrightarrow T^{1,0}X.
    \]
\end{lemma}

As a consequence of this, we derive a description of all first-order deformations of isomorphism classes of complex structures of $ X=(M,I) $.

\begin{proposition}
    Let $ X $ be a complex manifold. There is a natural bijection between all first-order deformations of $ X $ and elements of $ H^1(X,\mathcal{T}_X) $.
\end{proposition}

\begin{proof}
    Write $ X=(M,I) $. Then the first order deformations of $ I $ correspond to $ \bar{\partial} $-closed elements in $ \mathcal{A}^{0,1}(\mathcal{X}) $ and by the previous lemma those that define isomorphic ones differ by elements of the image of $ \bar{\partial}:\mathcal{A}^{0}(\mathcal{T}_X)\rightarrow\mathcal{A}^{0,1}(\mathcal{T}_X) $, i.e. $ \bar{\partial} $-exact ones. As the Dolbeault cohomology just computes the sheaf cohomology of $ \mathcal{T}_X $, this proves the assertion.
\end{proof}

In other words, the Kodaira-Spencer class determines the first-order deformations of $ X $ and any class in $ H^1(X,\mathcal{T}_X) $ occurs as a Kodaira-Spencer class.

The principal task in deformation theory is to integrate given first-order deformations $ v\in H^1(X,\mathcal{T}_X) $, i.e.\ to find a one-parameter family $ I_t $ such that its Kodaira-Spencer class is $ v $. In general this is not possible. Obstructions may occur at any order.

Moreover, since we have to divide by the action of the diffeomorphism group at any step, we have to study all possible $ \phi_i $ carefully. This is needed in order to ensure the convergence of the power series and to solve the recursive equations.

That there in general exist obstructions can already be guessed at order two: Once a lift $ \phi_1 $ of $ v\in H^1(X,\mathcal{T_X}) $ is chosen, we have to find $ \phi_2 $ such that $ \bar{\partial}\phi_2=-[\phi_1,\phi_1] $. But it may very well happen that the form $ [\phi_1,\phi_1] $ is never $ \bar{\partial} $-exact, no matter how $ \phi_1 $ is chosen.

Using~\ref{cvb} we can associate to any first order deformation $ v\in H^1(X,\mathcal{T}_X) $ a class $ [v,v]\in H^2(X,\mathcal{T}_X) $.

\begin{corollary}
    A first-order deformation $ v\in H^1(X,\mathcal{T}_X) $ cannot be integrated if $ [v,v]\in H^2(X,\mathcal{T}_X) $ does not vanish.
\end{corollary}

\subsection{Deformations of Calabi-Yau manifolds}

Next, we will describe how the recursive system of equations given by the Maurer-Cartan equation can be sovled when the underlying manifold $ X $ admits a holomorphic volume form, i.e.\ when $ X $ is a \textbf{Calabi-Yau manifold}. More precisely, by a Calabi-Yau manifold we mean here a compact K\"{a}hler manifold $ X $ of dimension $ n $ with trivial canonical bundle $ \bigwedge^n\varOmega_X=K_X \cong \mathcal{O}_X $. Moreover, we fix a holomorphic volume form, i.e.\ a trivializing section $ \varOmega\in H^0(X,K_X) $.

The form $ \varOmega $ can be used to define a natural isomorphism\[
    \eta:\bigwedge\nolimits^p\mathcal{T}_X\cong \bigwedge\nolimits^{n-p}\varOmega_X.
\]

One defines\[
    \eta(v_1\wedge\dots\wedge v_p)=i_{v_1}\dots i_{v_p}(\varOmega),
\]where $ i_v $ is the contraction defined by $ i_v(\alpha)(w_1,\dots,w_i)=\alpha(v,w_1,\dots,w_i) $. If $ \varOmega $ is locally written in the form $ fdz_1\wedge\dots\wedge dz_n $, then\[
    \eta(\frac{\partial}{\partial z_{i_1}}\wedge\dots\wedge\frac{\partial}{\partial z_{i_p}})={(-1)}^{(\sum i_j)-p}fdz_1\wedge\dots\wedge\widehat{dz_{i_1}}\wedge\dots\wedge\widehat{dz_{i_p}}\wedge\dots dz_n,
\]for $ i_1 < \cdots <i_p $.

Moreover, the isomorphism $ \eta $ induces canonical isomorphisms\[
    \eta:\mathcal{A}^{0,q}(\bigwedge\nolimits^p\mathcal{T}_X)\cong\mathcal{A}^{n-p,q}(X).
\]

\begin{definition}
    The operator $ \varDelta:\mathcal{A}^{0,q}(\bigwedge^p\mathcal{T}_X)\rightarrow \mathcal{A}^{0,q}(\bigwedge^{p-1}\mathcal{T}_X) $ is defined as\[
        \varDelta:\mathcal{A}^{0,q}(\bigwedge\nolimits^p\mathcal{T}_X)\stackrel{\eta}{\longrightarrow}\mathcal{A}^{n-p,q}(X)\stackrel{\partial}{\longrightarrow}\mathcal{A}^{n-p+1,q}(X)\stackrel{\eta^{-1}}{\longrightarrow}\mathcal{A}^{0,q}(\bigwedge\nolimits^{p-1}\mathcal{T}_X).
    \]
\end{definition}

Note that $ \partial $ itself is not well-defined on $ \mathcal{A}^{0,q}(\bigwedge^p\mathcal{T}_X) $.

\begin{lemma}
    The operator $ \varDelta $ anti-commutes with $ \bar{\partial} $, i.e.\ $ \varDelta\circ\bar{\partial}=-\bar{\partial}\circ\varDelta $.
\end{lemma}

\begin{proof}
    We use the fact that $ \varOmega\in\mathcal{A}^{n,0}(X) $ is holomorphic, i.e.\ $ \bar{\partial} $-closed. If we write in local coordinates $ \varOmega=fdz_1\wedge\dots\wedge dz_n $, then $ f $ is a holomorphic function. Then\[
        \eta(g\cdot d\bar{z}_{I}\otimes\frac{\partial}{\partial z_J})=\pm(fg)\cdot d\bar{z}_{I}\otimes dz_{\{1,\dots,n\}\backslash J}.
    \]Hence, $ \bar{\partial}\circ\eta=\eta\circ\bar{\partial} $ and, therefore,\begin{align*}
        \bar{\partial}(\varDelta(\alpha))&=\eta^{-1}(\bar{\partial}\partial\eta(\alpha))=-\eta^{-1}(\partial\bar{\partial}\eta(\alpha))\\
        &=-\eta^{-1}(\partial\eta(\bar{\partial}\alpha))=-\varDelta(\bar{\partial}(\alpha)).
    \end{align*}
\end{proof}

Also note that $ \varDelta^2=0 $. Thus, $ \varDelta $ can be used as a differential, but it is not well-behaved with respect to the exterior product $ \mathcal{A}^{0,p}(\bigwedge^r \mathcal{T}_X)\times\mathcal{A}^{0,q}(\bigwedge^s \mathcal{T}_X)\rightarrow\mathcal{A}^{0,p+q}(\bigwedge^{r+s}\mathcal{T}_X) $. This is made more precise by the following key result due to Tian and Todorov~\cite{Tian87,TT89}.

\begin{lemma}[Tian-Todorov lemma]
    If $ \alpha\in\mathcal{A}^{0,p}(\mathcal{T}_X) $ and $ \beta\in\mathcal{A}^{0,q}(\mathcal{T}_X) $, then\[
        {(-1)}^p [\alpha,\beta]=\varDelta(\alpha\wedge\beta)-\varDelta(\alpha)\wedge\beta-{(-1)}^{p+1} \alpha\wedge\varDelta(\beta).
    \]
\end{lemma}

\begin{proof}
    As the claim is local and additive in $ \alpha $ and $ \beta $, it suffices to prove it for $ \alpha=ad\bar{z}_I\otimes\frac{\partial}{\partial z_i} $ and $ \beta=bd\bar{z}_J\otimes\frac{\partial}{\partial z_j} $.

    Then, by definition, $ [\alpha,\beta]=d\bar{z}_I \wedge d\bar{z}_J [a\frac{\partial}{\partial z_i},b\frac{\partial}{\partial z_j}] $.
    
    In a first step we shall reduce the assertion to the case $ p=q=0 $. Using the short hand \[
        G(\alpha,\beta):=\varDelta(\alpha\wedge\beta)-\varDelta(\alpha)\wedge\beta-{(-1)}^{p+1}\alpha\wedge\varDelta(\beta),
    \]one does this by proving\[
        G(\alpha,\beta)=(-1)^p (d\bar{z}_I \wedge d\bar{z}_J)\otimes G(a\frac{\partial}{\partial z_i},b\frac{\partial}{\partial z_j}).
    \]Using the standard sign rules, one indeed computes\begin{align*}
        \varDelta(\alpha\beta)&=\eta^{-1}\partial\eta(\alpha\wedge\beta)=(-1)^q \eta^{-1}\partial\eta\bigg((d\bar{z}_I \wedge d\bar{z}_J)\otimes(a\frac{\partial}{\partial z_i}\wedge b\frac{\partial}{\partial z_j})\bigg)\\
        &=(-1)^q \eta^{-1}\partial\bigg((d\bar{z}_I \wedge d\bar{z}_J)\otimes \eta\bigg(a\frac{\partial}{\partial}\wedge b\frac{\partial}{\partial z_j}\bigg)\bigg)\\
        &=(-1)^q \cdot (-1)^{p+q} \eta^{-1}\bigg((d\bar{z}_I \wedge d\bar{z}_J)\otimes\partial\eta\bigg(a\frac{\partial}{\partial z_i}\wedge b\frac{\partial}{\partial z_j}\bigg)\bigg).\\
        &=(-1)^p (d\bar{z}_I \wedge d\bar{z}_J)\otimes\varDelta\bigg(a\frac{\partial}{\partial z_i}\wedge b\frac{\partial}{\partial z_j}\bigg).
    \end{align*}
    Similarly,\begin{align*}
        \varDelta(\alpha)\wedge\beta&=\eta^{-1}\partial\eta\bigg(d\bar{z}_I \otimes a\frac{\partial}{\partial z_i}\bigg)\wedge\beta=(-1)^p \bigg(d\bar{z}_I \otimes \varDelta\bigg( a\frac{\partial}{\partial z_i}\bigg)\bigg)\wedge\beta\\
        &=(-1)^p (d\bar{z}_I \wedge d\bar{z}_J)\otimes \bigg(\varDelta \bigg(a\frac{\partial}{\partial z_i}\bigg)\wedge b\frac{\partial}{\partial z_j}\bigg)
    \end{align*}and\begin{align*}
        (-1)^{p+1}\alpha\wedge\varDelta(\beta)&=(-1)^{p+1}\alpha\wedge\bigg((-1)^q d\bar{z}_J \otimes \varDelta\bigg(b\frac{\partial}{\partial z_j}\bigg)\bigg)\\
        &=(-1)^p (-1)(d\bar{z}_I \wedge d\bar{z}_J)\otimes \bigg(a\frac{\partial}{\partial  z_i}\wedge\varDelta(b\frac{\partial}{\partial z_j})\bigg).
    \end{align*}Thus,\begin{align*}
        G(\alpha,\beta)=&(-1)^p (d\bar{z}_I \wedge d\bar{z}_J)\otimes \bigg(\varDelta\bigg(a\frac{\partial}{\partial z_i}\wedge b\frac{\partial}{\partial z_j}\bigg)\\ &-\varDelta\bigg(a\frac{\partial}{\partial z_i}\bigg)\wedge b\frac{\partial}{\partial z_j}-(-1)a\frac{\partial}{\partial z_i}\wedge\varDelta\bigg(b\frac{\partial}{\partial z_j}\bigg)\bigg)\\
        =&(-1)^p (d\bar{z}_I \wedge d\bar{z}_J)\otimes G\bigg(a\frac{\partial}{\partial z_i},b\frac{\partial}{\partial z_j}\bigg).
    \end{align*}

    In order to treat the case $ p=q=0 $, let us first mention explicitly\[
        \eta(\frac{\partial}{\partial z_i})=(-1)^{i-1} fdz_1\wedge\dots\wedge\widehat{dz_i}\wedge\dots\wedge dz_n
    \]and\[
        \eta\bigg(\frac{\partial}{\partial z_i}\wedge\frac{\partial}{\partial z_j}\bigg)=(-1)^{i+j-2} fdz_1 \wedge \dots\wedge\widehat{dz_i}\wedge\dots\wedge\widehat{dz_j}\wedge\dots\wedge dz_n.
    \]This also shows $ \eta^{-1}(dz_1\wedge\dots\wedge\widehat{dz_i}\wedge\dots dz_n)=f^{-1}\frac{\partial}{\partial z_i} $ and $ \eta(dz_1\wedge\dots\wedge dz_n)=f^{-1} $.

    Hence,\begin{align*}
        \varDelta(a\frac{\partial}{\partial z_i})&=\eta^{-1}\partial\Big((-1)^{i-1}(af)dz_1\wedge\dots\wedge\widehat{dz_i}\wedge\dots\wedge dz_n\Big)\\
        &=\eta^{-1}\frac{\partial(af)}{\partial z_i}dz_1 \wedge \dots\wedge dz_n=\frac{\partial(af)}{\partial z_i}f^{-1}
    \end{align*}and similarly\[
        \varDelta\bigg(b\frac{\partial}{\partial z_j}\bigg)=\frac{\partial(bf)}{\partial z_j}f^{-1}.
    \]Furthermore,\begin{align*}
        \varDelta\bigg(a\frac{\partial}{\partial z_i},b\frac{\partial}{\partial z_j}\bigg)=&\eta^{-1}\partial\Big((-1)^{i+j}(abf)dz_1 \wedge\dots\wedge\widehat{dz_i}\wedge\dots\wedge\widehat{dz_j}\wedge\dots\wedge dz_n\Big)\\
        =&\eta^{-1}\bigg((-1)^{j-1}\frac{\partial(abf)}{\partial z_i}dz_1 \wedge\dots\wedge\widehat{dz_j}\wedge\dots\wedge dz_n\\
        &+(-1)^i \frac{\partial(abf)}{\partial z_j}dz_1\wedge\dots\wedge\widehat{dz_i}\wedge\dots\wedge dz_n\bigg)\\
        =&\frac{\partial(abf)}{\partial z_i}f^{-1}\frac{\partial}{\partial z_j}-\frac{\partial(abf)}{\partial z_j}f^{-1}\frac{\partial}{\partial z_i}\\
        =&\bigg[a\frac{\partial}{\partial z_i},b\frac{\partial}{\partial z_j}\bigg]+\varDelta\bigg(a \frac{\partial}{\partial z_i}\bigg)b\frac{\partial}{\partial z_j}-a\frac{\partial}{\partial z_i}\varDelta\bigg(b\frac{\partial}{\partial z_j}\bigg).
    \end{align*}This proves the assertion for $ p=q=0 $.
\end{proof}

\begin{corollary}
    Let $ \alpha $ and $ \beta $ be $ \varDelta $-closed elements of $ \mathcal{A}^{0,*}(\mathcal{T}_X) $. Then the bracket $ [\alpha,\beta] $ is $ \varDelta $-exact.
\end{corollary}

The Tian-Todorov lemma will be applied to the coefficients $ \phi_i\in \mathcal{A}^{0,1}(\mathcal{T}_X) $ of the power series expansion of a deformation of the complex structure. Before proving the general result, we shall spell out the first two steps.

Let $ v\in H^1(X,\mathcal{T}_X) $ be given. We first look for a good representative $ \phi_1\in\mathcal{A}^{0,1}(\mathcal{T}_X) $. In order to solve the Maurer-Cartan equation $ \bar{\partial}\phi_2=-[\phi_1,\phi_1] $, we have to ensure that\begin{itemize}
    \item $ [\phi_1,\phi_1] $ is $ \bar{\partial} $-closed
    \item $ [\phi_1,\phi_1] $ has no harmonic part.
\end{itemize}

Here we use the Hodge decomposition $ \mathcal{A}^{0,1}(\mathcal{T}_X)=\bar{\partial}\mathcal{A}^{0}(\mathcal{T}_X)\oplus\mathcal{H}^{0,1}(\mathcal{T}_X)\oplus\bar{\partial}^*\mathcal{A}^{0,2}(\mathcal{T}_X) $ with respect to an hermitian metric on $ X $ and thus on $ \mathcal{T}_X $.

The first condition is always satisfied: If $ \alpha_i\in\mathcal{A}^{0,q_i}(\mathcal{T}_X),\, i=1,2 $, are $ \bar{\partial} $-closed, then so is $ [\alpha_1,\alpha_2] $. This follows from (\ref{cvb}).

The Tian-Todorov lemma yields: If $ \alpha_i\in\mathcal{A}^{0,q_i}(\mathcal{T}_X),\, i=1,2 $ are chosen such that $ \eta(\alpha_i)\in\mathcal{A}^{n-1,q_i}(X) $ are $ \partial $-closed, then $ \eta[\alpha_1,\alpha_2] $ is $ \partial $-exact. More explicitly, $ \eta[\alpha_1,\alpha_2]=(-1)^{q_1}\cdot \partial\eta(\alpha_1 \wedge\alpha_2) $.

Eventually, we find: If $ \phi_1\in \mathcal{A}^{0,1}(\mathcal{T}_X) $ is chosen such that $ \eta(\phi_1)\in\mathcal{A}^{n-1,1}(X) $ is harmonic with respect to a K\"{a}hler metric on $ X $, then $ \bar{\partial}\phi_2=-[\phi_1,\phi_1] $ admits a solution $ \phi_2\in\mathcal{A}^{0,1}(\mathcal{T}_X) $ which can, moreover, be chosen such that $ \eta(\phi_2) $ is $ \partial $-exact.

Indeed, by the above arguments $ \eta[\phi_1,\phi_1] $ is $ \partial $-exact and $ \bar{\partial} $-closed. Since by Hodge theory on a compact K\"{a}hler manifold a $ \partial $-exact form cannot have a non-trivial harmonic part, this proves that $ \eta[\phi_1,\phi_1] $ is $ \bar{\partial} $-exact. But using the $\partial\bar{\partial} $-lemma, we can, moreover, say that there exists a form $ \gamma\in\mathcal{A}^{n-2,1}(X) $ with $ \bar{\partial}\partial\gamma=-\eta[\phi_1,\phi_1] $. Then set $ \phi_2:=\eta^{-1}(\partial\gamma) $. Since $ \eta $ commutes with $ \bar{\partial} $, this proves the claim.

All these arguments can now easily be generalized to prove the general result, which is also due to Tian and Todorov.

\begin{proposition}
    Let $ X $ be a Calabi-Yau manifold and let $ v\in H^1(X,\mathcal{T}_X) $. Then there exists a formal power series $ \phi_1 t+\phi_2 t^2+\cdots $ with $ \phi_i\in\mathcal{A}^{0,1}(\mathcal{T}_X) $ satisfying the Maurer-Cartan equations\[
        \bar{\partial}\phi_1=0\ \text{and}\ \bar{\partial}\phi_k=-\sum_{0<i<k}[\phi_i,\phi_{k-i}],
    \]with $ [\phi_1]=v $ and such that $ \eta(\phi_i)\in\mathcal{A}^{n-1,1}(X) $ is $ \partial $-exact for all $ i>1 $.
\end{proposition}

\begin{proof}
    We begin with $ \phi_1\in\mathcal{A}^{0,1}(\mathcal{T}_X) $ representing $ v $ such that $ \eta(\phi_1)\in\mathcal{A}^{n-1,1}(X) $ is harmonic. Then one finds $ \phi_2 $ as above. Let us suppose that we have found $ \phi_2,\dots,\phi_{k-1} $ as claimed by the proposition. Then $ \eta[\phi_i,\phi_{k-i}] $ is $ \partial $-exact for $ 0<i<k $ by the Tian-Todorov lemma. Thus, in order to find $ \phi_k $ satisfying the assertion, we have to ensure that the sum $ \sum_{0<i<k}[\phi_i,\phi_{k-i}] $ is $ \bar{\partial} $-closed. As we have used implicitly before, one has\[
        \bar{\partial}\bigg(\sum_{0<i<k}[\phi_i,\phi_{k-i}]\bigg)=\sum_{0<i<k}([\bar{\partial}\phi_i,\phi_{k-i}]-[\phi_{i},\bar{\partial}\phi_{k-i}])
    \]and by induction hypothesis this equals\begin{align*}
        & -\sum_{0<i<k}\bigg(\sum_{0<j<i}[[\phi_j,\phi_{i-j}],\phi_{k-i}]-\sum_{0<l<k-i}[\phi_i,[\phi_l,\phi_{k-i-l}]]\bigg)\\
        =&-\sum_{0<i<k}\sum_{0<j<i}[[\phi_j,\phi_{i-j}],\phi_{k-i}]+\sum_{0<i<k}\sum_{0<l<i}[\phi_{k-i},[\phi_l,\phi_{i-l}]]\\
        =&2\sum_{0<i<k}\sum_{0<l<i}[\phi_{k-i},[\phi_l,\phi_{i-l}]]=0
    \end{align*}where we use $ [\alpha,\beta]=-[\beta,\alpha] $ for $ \alpha\in\mathcal{A}^{0,2}(\mathcal{T}_X) $ and arbitrary $ \beta $ and by Jacobi identity we deduce that the final result is 0.

    The $ \partial\bar{\partial} $-lemma eventually ensures that we can choose $ \phi_k $ of the form $ \eta^{-1}(\partial\gamma) $, i.e.\ $ \eta(\phi_k) $ is $ \partial $-exact.
\end{proof}

\begin{remark}
    There might be very well other solutions of the Maurer-Cartan equation which do not satisfy the extra condition that $ \eta(\phi_i) $ is $ \partial $-exact for $ i>1 $. In fact, the $ \eta(\phi_i) $ need not even be $ \partial $-closed.
\end{remark}

\begin{remark}
    So far, we have not said anything about the existence of a convergent solution. In fact, there is a standard procedure to turn any formal solution into a convergent one. At this point, a good deal of analysis arises. For the details we refer to~\cite[Sect. 5.3]{Kodaira}.

    The principal idea is to consider the operator $ \bar{\partial} $ on $ \mathcal{A}^{0,q}(\mathcal{T}_X) $ and its adjoint operator with respect to a chosen K\"{a}hler (or just hermitian) metric. For clarity, let us write $ \bar{\partial}_{\mathcal{T}_X} $ respectively $ \bar{\partial}^*_{\mathcal{T}_X} $ for these two operators. Then one shows that a formal solution converges if the coefficients $ \phi_k $ are $ \bar{\partial}^*_{\mathcal{T}_X} $-exact. It is then not too difficult to transform any formal solution into one that satisfies this additional condition, as the mere existence of a formal solution ensures that all possible obstructions are trivial.
\end{remark}

\section{Deformation theory with algebraic approaches}

\subsection{The $ T^i $ functors}

We first present the construction and main properties of the $ T^i $ functors introduced by Lichtenbaum and Schlessinger~\cite{LichtenbaumSchlessinger}.

Let $ A\rightarrow B $ be a homomorphism of rings and let $ M $ be a $ B $-module. Here we construct groups $ T^i(B/A,M) $ for $ i=0,1,2 $. The rings are assumed to be commutative with identity, but we do not impose any finiteness conditions yet.

First choose a polynomial ring $ R=A[x] $ in a set of variables $ x=\{x_i\} $ (possibly infinite) such that $ B $ can be written as a quotient of $ R $ as an $ A $-algebra. Let $ I $ be the ideal defining $ B $, so that there is an exact sequence\[
    0\longrightarrow I\longrightarrow R\longrightarrow B\longrightarrow 0.
\]Second choose a free $ R $-module $ F $ and a surjection $ j:F\rightarrow I\rightarrow 0 $ and let $ Q $ be the kernel:\[
    0\longrightarrow Q\longrightarrow F\stackrel{j}{\longrightarrow}I\longrightarrow 0.
\]Let\[
    F_0=(j(a)b-j(b)a)_{a,b\in F}\subseteq F.
\]Define a complex of $ B $-modules, called the \textbf{cotangent complex}\[
    L_2\stackrel{d_2}{\longrightarrow}L_1\stackrel{d_1}{\longrightarrow}L_0
\]as follows.\begin{itemize}
    \item Take $ L_2=Q/F_0 $, which is a $ B $-module. Indeed, if $ x\in I $ and $ a\in Q $. Write $ x=j(x') $ for some $ x'\in F $, then\[
        xa=j(x')a\equiv j(a)x'=0\pmod{F_0}.
    \]
    \item Take $ L_1=F\otimes_R B=F/IF $, and let $ d_2:L_2\rightarrow L_1 $ be the map induced from the inclusion $ Q\rightarrow F $.
    \item Take $ L_0=\varOmega_{R/A}\times_R B $, where $ \varOmega_{R/A} $ is the module of relative K\"{a}hler differential. Let $ d_1 $ be the composition of\begin{align*}
        L_1\twoheadrightarrow I/I^2 &\rightarrow \varOmega_{R/A}\otimes_R B\\
        x &\rightarrow d_{R/A}(x)\otimes 1
    \end{align*}It is clear that $ d_1 d_2=0 $, so we defined a complex of $ B $-modules. Notice that $ L_0 $, $ L_1 $ are free $ B $-modules, since $ R $ is a polynomial ring over $ A $.
    \item Indeed, the second fundamental exact sequence of K\"{a}hler differential\[
        I/I^2\stackrel{d}{\longrightarrow}\varOmega_{R/A}\otimes_R B\longrightarrow \varOmega_{B/A}\longrightarrow 0
    \]give the exact sequence\begin{equation}\label{eq10}
        L_2\stackrel{d_2}{\longrightarrow}L_1\stackrel{d_1}{\longrightarrow}L_0\longrightarrow\varOmega_{B/A}\longrightarrow 0.
    \end{equation}
\end{itemize}

\begin{lemma}\label{A1}
    For any $ B $-module $ M $, define\[
        T^i(B/A,M):=H^i(\mathrm{Hom}_B(L_{\bullet},M)),
    \]which is independent of the choice of $ F $ and $ R $.
\end{lemma}

\begin{proof}
    First, we fixed the polynomial $ R $. If $ F $ and $ F'
     $ are two choices of free $ R $-modules mapping onto $ I $, then $ F\oplus F' $ gives the third choice. Since $ F' $ is free, the map $ j':F'\rightarrow I $ factors through $ F $, i.e., $ j'=jp $ for some map $ p:F'\rightarrow F $. Replacing each generator $ e' $ of $ F'\subseteq F\oplus F' $ by $ e'-p(e') $, we may assume that the map $ F\oplus F'\rightarrow I $ is just $ (j,0) $. Then $ \ker(j,0)=Q\oplus F' $ and thus $ (F \oplus F')_0=F_0+IF' $. Let \(L_{\bullet}^{\prime}\) be the complex obtained from \((F\oplus F^{\prime},R)\), then\[
        L_{2}^{\prime} = L_{2}\oplus F^{\prime} / I F^{\prime},\ L_{1}^{\prime} = L_{1}\oplus (F^{\prime}\otimes_{R}B),\ L_{0}^{\prime} = L_{0}.
    \]
    Hence, \(L_{\bullet}^{\prime}\) is obtained by the direct sum of \(L_{\bullet}\) with the free acyclic complex \(F^{\prime}\otimes_{R}B\to F^{\prime}\otimes_{R}B\to 0\), which induces same cohomology after taking $ \mathrm{Hom} $ functor.

    Second, it suffices to compare \(R = A[x]\) with \(R^{\prime} = A[x,y]\). Again, the map \(A[y]\to B\) can be factored through \(A[x]\) by a homomorphism \(p:A[y]\to A[x]\). Replacing each \(y_{i}\in A[x,y]\) by \(y_{i} - p(y_{i})\), we may assume the ring homomorphism \(A[x,y]\to B\) maps \(y_{i}\) to 0, then it has kernel \(I R^{\prime} + y R^{\prime}\). Since \(T^{i}\) is independent of the choice of \(F\), we may take \(F^{\prime} = F\otimes_{R}R^{\prime}\) and \(G\) be a free \(R^{\prime}\)-module on the index set of the \(y\) variables. Then we have the following diagram% https://quiver.taskfirst.cn/#q=WzAsMTAsWzAsMCwiMCJdLFsxLDAsIlEnIl0sWzIsMCwiRidcXG9wbHVzIEcnIl0sWzMsMCwiSVInXFxvcGx1cyB5UiciXSxbNCwwLCIwIl0sWzQsMSwiMCJdLFszLDEsIkkiXSxbMiwxLCJGIl0sWzEsMSwiUSJdLFswLDEsIjAiXSxbMCwxXSxbMSwyXSxbOSw4XSxbOCwxXSxbOCw3XSxbNywyXSxbNyw2XSxbNiwzXSxbNiw1XSxbMyw0XSxbMiwzXV0=
    \[\begin{tikzcd}
        0 & {Q'} & {F'\oplus G'} & {IR'\oplus yR'} & 0 \\
        0 & Q & F & I & 0
        \arrow[from=1-1, to=1-2]
        \arrow[from=1-2, to=1-3]
        \arrow[from=1-3, to=1-4]
        \arrow[from=1-4, to=1-5]
        \arrow[from=2-1, to=2-2]
        \arrow[from=2-2, to=1-2]
        \arrow[from=2-2, to=2-3]
        \arrow[from=2-3, to=1-3]
        \arrow[from=2-3, to=2-4]
        \arrow[from=2-4, to=1-4]
        \arrow[from=2-4, to=2-5]
    \end{tikzcd}\]and it is clear that\[
        Q^{\prime} = Q R^{\prime} + \langle y_{i}a - j(a)e_{i}\mid a\in F\rangle_{R^{\prime}} + \langle y_{i}e_{j} - y_{j}e_{i}\rangle_{R^{\prime}}.
    \]
    Since \(\langle y_{i}a - j(a)e_{i}\mid a\in F\rangle_{R^{\prime}} + \langle y_{i}e_{j} - y_{j}e_{i}\rangle_{R^{\prime}}\subseteq (F^{\prime}\oplus G^{\prime})_{0}\), we have \(L_{2}^{\prime} = L_{2}\).

    On the other hand, \(L_{1}^{\prime} = L_{1}\oplus (G^{\prime}\otimes_{R^{\prime}}B)\) and \(L_{0}^{\prime} = L_{0}\oplus (\varOmega_{A[y] / A}\otimes B)\). Then \(L_{\bullet}^{\prime}\) is obtained by the direct sum of \(L_{\bullet}\) with the free acyclic complex\begin{align*}
        G^{\prime}\otimes_{R^{\prime}}B&\longrightarrow \varOmega_{A[y] / A}\otimes B\longrightarrow 0\\
        e_{i}&\longrightarrow dy_{i},
    \end{align*}
    and thus induces same cohomology after taking $ \mathrm{Hom} $ functor.
\end{proof}

\begin{theorem}
Let \(A \to B\) be a homomorphism of rings. Then \(T^{i}(B / A, - )\) is a covariant, additive functor from \(\mathsf{Mod}_{B}\) to itself. If
\[
0\longrightarrow M^{\prime}\longrightarrow M\longrightarrow M^{\prime \prime}\longrightarrow 0
\]
is a short exact sequence of \(B\)-modules, then there is a long exact sequence
\[
\begin{aligned}
0 &\longrightarrow T^{0}(B / A,M^{\prime})\longrightarrow T^{0}(B / A,M)\longrightarrow T^{0}(B / A,M^{\prime \prime}) \\
&\longrightarrow T^{1}(B / A,M^{\prime})\longrightarrow T^{1}(B / A,M)\longrightarrow T^{1}(B / A,M^{\prime \prime}) \\
&\longrightarrow T^{2}(B / A,M^{\prime})\longrightarrow T^{2}(B / A,M)\longrightarrow T^{2}(B / A,M^{\prime \prime}).
\end{aligned}
\]
\end{theorem}

\begin{proof}
By construction, \(T^{i}(B / A, - )\) is a covariant additive functor. Since the terms \(L_{1}\) and \(L_{0}\) of the complex \(L_{\bullet}\) are free, we get a sequence of complexes
\[
0\longrightarrow \mathrm{Hom}_{B}(L_{\bullet},M^{\prime})\longrightarrow \mathrm{Hom}_{B}(L_{\bullet},M)\longrightarrow \mathrm{Hom}_{B}(L_{\bullet},M^{\prime \prime})\longrightarrow 0
\]
that is exact except possibly for the map
\[
\mathrm{Hom}_{B}(L_{2},M)\to \mathrm{Hom}_{B}(L_{2},M^{\prime \prime})
\]
may not be surjective. This sequence of complexes gives the long exact sequence of cohomology above.
\end{proof}

\begin{theorem}
    Let \(A \to B \to C\) be homomorphisms of rings, and let \(M\) be a \(C\)-module. Then there is a long exact sequence\[
        \begin{aligned}
            0 &\longrightarrow T^{0}(C / B,M)\longrightarrow T^{0}(C / A,M)\longrightarrow T^{0}(B / A,M) \\
            &\longrightarrow T^{1}(C / B,M)\longrightarrow T^{1}(C / A,M)\longrightarrow T^{1}(B / A,M) \\
            &\longrightarrow T^{2}(C / B,M)\longrightarrow T^{2}(C / A,M)\longrightarrow T^{2}(B / A,M).
        \end{aligned}
    \]
\end{theorem}

\begin{proof}
    Choose\[
        0\longrightarrow I\longrightarrow A[x]\longrightarrow B\longrightarrow 0,\quad 
        0\longrightarrow Q\longrightarrow F\stackrel{j}{\longrightarrow} I\longrightarrow 0
    \]\[
        0\longrightarrow J\longrightarrow B[y]\longrightarrow C\longrightarrow 0,\quad 
        0\longrightarrow P\longrightarrow G\longrightarrow J\longrightarrow 0
    \]to calculate \(T^{i}(B / A, M)\) and \(T^{i}(C / B, M)\). Then \(A[x, y] \to B[y] \to C\) is surjective, and denote its kernel by \(K\). It induces an exact sequence\[
        0\longrightarrow I[y]\longrightarrow K\longrightarrow J\longrightarrow 0
    \]of \(A\)-modules. Let \(F^{\prime}\) and \(G^{\prime}\) be free \(A[x,y]\)-modules on the same index sets as \(F\) and \(G\) respectively. Choose a lifting of the map \(G\to J\) to a map \(G^{\prime}\to K\). From the map \(F\to I\) we get the natural map \(F^{\prime}\to I[y]\). Hence, we get the surjection \(F^{\prime}\oplus G^{\prime}\to K\), and let \(S\) be its kernel. Now, we can calculate \(T^{i}(C/A, M)\) by\[
        0 \longrightarrow K \longrightarrow A[x,y] \longrightarrow C \longrightarrow 0,\quad 0 \longrightarrow S \longrightarrow F^{\prime}\oplus G^{\prime} \stackrel{j}{\longrightarrow} K \longrightarrow 0.
    \]Now, we have the following diagram% https://quiver.taskfirst.cn/#q=WzAsMTUsWzAsMCwiMCJdLFswLDEsIjAiXSxbMCwyLCIwIl0sWzEsMCwiUVt5XSJdLFsyLDAsIkYnIl0sWzMsMCwiSVt5XSJdLFs0LDAsIjAiXSxbNCwxLCIwIl0sWzQsMiwiMCJdLFsxLDIsIlAiXSxbMiwyLCJHIl0sWzMsMiwiSiJdLFszLDEsIksiXSxbMiwxLCJGJ1xcb3BsdXMgRyciXSxbMSwxLCJTIl0sWzAsM10sWzMsNF0sWzQsNV0sWzUsNl0sWzEyLDddLFsxLDE0XSxbMTQsMTNdLFsxMywxMl0sWzMsMTRdLFs0LDEzXSxbNSwxMl0sWzE0LDldLFsxMywxMF0sWzEyLDExXSxbMiw5XSxbOSwxMF0sWzEwLDExXSxbMTEsOF1d
    \[\begin{tikzcd}
        0 & {Q[y]} & {F'} & {I[y]} & 0 \\
        0 & S & {F'\oplus G'} & K & 0 \\
        0 & P & G & J & 0
        \arrow[from=1-1, to=1-2]
        \arrow[from=1-2, to=1-3]
        \arrow[from=1-2, to=2-2]
        \arrow[from=1-3, to=1-4]
        \arrow[from=1-3, to=2-3]
        \arrow[from=1-4, to=1-5]
        \arrow[from=1-4, to=2-4]
        \arrow[from=2-1, to=2-2]
        \arrow[from=2-2, to=2-3]
        \arrow[from=2-2, to=3-2]
        \arrow[from=2-3, to=2-4]
        \arrow[from=2-3, to=3-3]
        \arrow[from=2-4, to=2-5]
        \arrow[from=2-4, to=3-4]
        \arrow[from=3-1, to=3-2]
        \arrow[from=3-2, to=3-3]
        \arrow[from=3-3, to=3-4]
        \arrow[from=3-4, to=3-5]
    \end{tikzcd}\]which is exact for each horizontal sequence, but not exact for vertical sequence. It induces the maps of complexes\[
        L_{\bullet}(B/A) \otimes_B C \to L_{\bullet}(C/A) \to L_{\bullet}(C/B).
    \]Indeed, on the degree 0 level we have split exact sequence.% https://quiver.taskfirst.cn/#q=WzAsNyxbMCwwLCJcXE9tZWdhX3tBW3hdL0F9XFxvdGltZXNfe0FbeF19IEMiXSxbMSwwLCJcXE9tZWdhX3tBW3gseV0vQX1cXG90aW1lc197QVt4LHldfSBDIl0sWzIsMCwiXFxPbWVnYV97Qlt5XS9CfVxcb3RpbWVzX3tCW3ldfSBDIl0sWzAsMSwiZHhfaSJdLFsxLDEsImR4X2kiXSxbMSwyLCJkeF9pLFxcIGR5X2oiXSxbMiwyLCIwLFxcIGR5X2oiXSxbMCwxXSxbMSwyXSxbMyw0LCIiLDEseyJzdHlsZSI6eyJ0YWlsIjp7Im5hbWUiOiJtYXBzIHRvIn19fV0sWzUsNiwiIiwxLHsic3R5bGUiOnsidGFpbCI6eyJuYW1lIjoibWFwcyB0byJ9fX1dXQ==
    \[\begin{tikzcd}
        {\varOmega_{A[x]/A}\otimes_{A[x]} C} & {\varOmega_{A[x,y]/A}\otimes_{A[x,y]} C} & {\varOmega_{B[y]/B}\otimes_{B[y]} C} \\
        {dx_i} & {dx_i} \\
        & {dx_i,\ dy_j} & {0,\ dy_j}
        \arrow[from=1-1, to=1-2]
        \arrow[from=1-2, to=1-3]
        \arrow[maps to, from=2-1, to=2-2]
        \arrow[maps to, from=3-2, to=3-3]
    \end{tikzcd}\]On the degree 1 level we have split exact sequence\[
        F \otimes_{A[x]} C \longrightarrow (F^{\prime}\oplus G^{\prime})\otimes_{A[x,y]} C \longrightarrow G \otimes_{B[y]} C
    \]of free \(C\)-modules. On the degree 2 level we have\begin{equation}\label{xbm}
        (Q/F_0)\otimes_B C \longrightarrow S/(F^{\prime}\oplus G^{\prime})_0 \longrightarrow P/G_0.
    \end{equation}We claim that \(S\to P\) is surjective. Indeed, if \(p\in P\). Take \((f^{\prime}, g^{\prime})\in F^{\prime}\oplus G^{\prime}\) mapping to \(p\). Then the image of \((f^{\prime}, g^{\prime})\) in \(K\) is contained in \(I[y]\). Since \(F^{\prime}\to I[y]\) is surjective, there exists \(\tilde f\in F^{\prime}\) having the same image in \(K\) as \((f^{\prime}, g^{\prime})\). Hence, \((f^{\prime}-\tilde f, g^{\prime})\in S\) maps to \(p\).

    To complete the proof, we remain to show the exactness of (\ref{xbm}) in the middle. Let \(s = f + g\) be an element of \(S\subseteq F^{\prime}\oplus G^{\prime}\) such that its image in \(P\) is contained in \(G_0\). Then the image \(g\) of \(g^{\prime}\) in \(G\) belongs to \(G_0\), write \(g = j(a)b - j(b)a\) for some \(a,b\in G\). Lift \(a,b\) to elements \(a^{\prime},b^{\prime}\) in \(G^{\prime}\), then \(j(a^{\prime})b^{\prime} - j(b^{\prime})a^{\prime}\in S\). Let \(g^{\prime\prime} = g - j(a^{\prime})b^{\prime} + j(b^{\prime})a^{\prime}\), then \(g^{\prime\prime}\in \ker(G^{\prime}\to G) = IG^{\prime}\). Say \(g^{\prime\prime} = xh\) and \(x = j(x^{\prime})\) for some \(x\in I\), \(h\in G^{\prime}\), and \(x^{\prime}\in F\). Then
    \[
    xh = j(x^{\prime})h \equiv j(h)x^{\prime} \pmod{(F^{\prime}+G^{\prime})_0}.
    \]
    Therefore,
    \[
    s \equiv f^{\prime} + g^{\prime\prime} \equiv f^{\prime} + j(h)x^{\prime} \pmod{(F^{\prime}+G^{\prime})_0}
    \]
    and \(f^{\prime} + j(h)x^{\prime}\in F^{\prime}\cap S = Q[y]\).
\end{proof}

\begin{proposition}\label{A5}
    For any morphism \(A \to B\) and \(M \in \mathsf{Mod}_{B}\),\[
        T^{0}(B / A,M) = \mathrm{Hom}_{B}(\varOmega_{B / A},M) = \mathrm{Der}_{A}(B,M).
    \]In particular,\[
        T^{0}(B / A,B) = \mathrm{Hom}_{B}(\varOmega_{B / A},B) = T_{B / A}
    \]is the tangent module of \(B\) over \(A\).
\end{proposition}

\begin{proof}
    Taking \(\mathrm{Hom}(-, M)\) to the exact sequence (\ref{eq10}), we have the exact sequence\[
        0\to \mathrm{Hom}(\varOmega_{B / A},M)\to \mathrm{Hom}(L_{0},M)\to \mathrm{Hom}(L_{1},M).
    \]
\end{proof}

\begin{proposition}\label{A6}
    If \(B\) is a polynomial ring over \(A\), then \(T^{i}(B / A, M) = 0\) for \(i = 1, 2\) and for all \(M\).
\end{proposition}

\begin{proof}
    We may take \(R = B\) and \(F = 0\) in our construction. Hence, the cotangent complex only has \(L_{0}\) term.
\end{proof}

\begin{proposition}\label{A7}
    If \(A \to B\) is a surjective ring homomorphism with kernel \(I\), then\[
        T^{0}(B / A,M) = 0,\ T^{1}(B / A,M) = \mathrm{Hom}_{B}(I / I^{2},M)
    \]for any \(M \in \mathsf{Mod}_{B}\). In particular, \(T^{1}(B / A, B) = \mathrm{Hom}_{B}(I / I^{2}, B)\) is the normal module \(N_{B / A}\) of \(\mathrm{Spec}\, B\) in \(\mathrm{Spec}\, A\).
\end{proposition}

\begin{proof}
    In this case, we can take \(R = A\), and thus \(L_{0} = 0\). Hence, \(T^{0} = 0\) for any module \(M\). Furthermore, the exact sequence\[
        0\longrightarrow Q\longrightarrow F\stackrel{j}{\longrightarrow}I\longrightarrow 0
    \]tensored with \(B\) gives% https://quiver.taskfirst.cn/#q=WzAsNixbMCwwLCJRXFxvdGltZXNfe0F9QiJdLFsxLDAsIkZcXG90aW1lc197QX1CIl0sWzIsMCwiSSAvIEleezJ9Il0sWzMsMCwiMCJdLFswLDEsIlEvRl8wIl0sWzEsMSwiRi9JRiJdLFswLDFdLFsxLDIsImoiXSxbMiwzXSxbNCw1XSxbMCw0LCIiLDAseyJzdHlsZSI6eyJoZWFkIjp7Im5hbWUiOiJlcGkifX19XSxbMSw1LCJcXHNpbSJdXQ==
    \[\begin{tikzcd}
        {Q\otimes_{A}B} & {F\otimes_{A}B} & {I / I^{2}} & 0 \\
        {Q/F_0} & {F/IF}
        \arrow[from=1-1, to=1-2]
        \arrow[two heads, from=1-1, to=2-1]
        \arrow["j", from=1-2, to=1-3]
        \arrow["\sim", from=1-2, to=2-2]
        \arrow[from=1-3, to=1-4]
        \arrow[from=2-1, to=2-2]
    \end{tikzcd}\]Hence, we have an exact sequence\[
        L_{2}\longrightarrow L_{1}\longrightarrow I / I^{2}\longrightarrow 0,
    \]and thus \(T^{1}(B / A, M) = \mathrm{Hom}_{B}(I / I^{2}, M)\).
\end{proof}

\begin{corollary}\label{A8}
    Suppose that \(A\) is a local ring and \(B = A / I\), where \(I\) is generated by a regular sequence \(a_{1},\ldots ,a_{r}\). Then \(T^{2}(B / A,M) = 0\) for all \(M\).
\end{corollary}

\begin{proof}
    Since the Koszul complex of a regular sequence is exact, we conclude that \(Q = F_{0}\). Hence, \(L_{2} = 0\) and \(T^{2}(B / A,M) = 0\) for all \(M\).
\end{proof}

\begin{proposition}\label{A9}
    Suppose that \(A = k[x_{1},\ldots ,x_{n}]\) and \(B = A / I\). Then for any \(M\), there is an exact sequence\[
        0\longrightarrow T^{0}(B / k,M)\longrightarrow \mathrm{Hom}_{A}(\varOmega_{A / k},M)\longrightarrow \mathrm{Hom}_{B}(I / I^{2},M)\longrightarrow T^{1}(B / k,M)\longrightarrow 0
    \]and an isomorphism\[
        T^{2}(B / A,M)\stackrel{\sim}{\longrightarrow} T^{2}(B / k,M).
    \]
\end{proposition}

\begin{proof}
    It follows from the long exact sequence of \(T^{i}\)-functors for \(k\to A\to B\) and (\ref{A5},~\ref{A6},~\ref{A7}).
\end{proof}

\begin{remark}\label{A10}
    We have not made any finiteness assumption on the rings and modules. It is clear that if $ A $ is a noetherian ring, $ B $ is a finitely generated $ A $-algebra, and $ M $ is a finitely generated $ B $-module, then $ T^i(B/A,M) $ are finitely generated $ B $-modules.
\end{remark}

\begin{lemma}\label{A11}
    The construction of the \(T^{i}\) functors is compatible with localization, and thus we may define sheaves \(\mathcal{T}^{i}(X / Y,\mathcal{F})\) for any morphism of schemes \(f:X\to Y\) and any sheaf \(\mathcal{F}\) of \(\mathcal{O}_{X}\)-modules, such that for any open affine \(\mathrm{Spec}\, A\subseteq Y\) and any open affine \(\mathrm{Spec}\, B\subseteq f^{-1}(\mathrm{Spec}\, A)\), where \(\mathcal{F}|_{\mathrm{Spec}\, B} = \widetilde{M}\), the section \(\mathcal{T}^{i}(X / Y,\mathcal{F})(\mathrm{Spec}\, B) = T^{i}(B / A,M)\).
\end{lemma}

For convenience, denote the modules \(T^{i}(B / A,B)\) and \(T^{i}(B / k,B)\) by \(T_{B / A}^{i}\) and \(T_{B / k}^{i}\) (or \(T_{B}^{i}\)) if no confusion arises. Furthermore, \(T_{B / A}^{0}\) will be written as the tangent module \(T_{B / A}\) of \(B\) over \(A\). Similarly for the sheaves \(\mathcal{T}^{i}(X / Y,\mathcal{O}_{X})\) and \(\mathcal{T}^{i}(X / k,\mathcal{O}_{X})\) we write \(\mathcal{T}_{X / Y}^{i}\) and \(\mathcal{T}_{X / k}^{i}\) (or \(\mathcal{T}_{X}^{i}\)). The sheaf \(\mathcal{T}_{X}^{0}\) will be written as the tangent sheaf \(\mathcal{T}_{X}\) of \(X\).

\begin{proposition}[Base change]\label{A12}
    Assume that \(A\) is Noetherian and \(B\) is a finitely generated \(A\)-algebra.\begin{enumerate}[(1)]
        \item Let \(M\in \mathsf{Mod}_{B}\) and \(A\to A^{\prime}\) be a flat morphism. Let \(B^{\prime} = B\otimes_{A}A^{\prime}\) and \(M^{\prime} = M\otimes_{B}B^{\prime}\). Then \(T^{i}(B / A,M)\otimes_{A}A^{\prime}\simeq T^{i}(B^{\prime} / A^{\prime},M^{\prime})\) for each \(i\).
        \item Let \(A \to A^{\prime}\) be an arbitrary ring homomorphism and \(B\) is flat over \(A\). Let \(B^{\prime} = B\otimes_{A}A^{\prime}\) and \(M^{\prime}\in \mathsf{Mod}_{B^{\prime}}\). Then \(T^{i}(B / A, M^{\prime}) = T^{i}(B^{\prime} / A^{\prime}, M^{\prime})\) for each \(i\).
    \end{enumerate}
\end{proposition}

\begin{proof}
    \cite[Ex. 3.7, 3.8]{Har10}
\end{proof}

An \textbf{infinitesimal thickening} $ Y\subseteq Y' $ means that $ Y $ is a closed subscheme of another scheme $ Y' $, and that the ideal $ \mathcal{I} $ defining $ Y $ inside $ Y' $ is nilpotent.

\begin{proposition}[Infinitesimal Lifting Property]\label{A13}
    Let \(X\) be a nonsingular affine scheme of finite type over \(k\), let \(f: Y \to X\) be a morphism from an affine scheme \(Y\) over \(k\), and let \(Y \subseteq Y'\) be an infinitesimal thickening of \(Y\). Then the morphism \(f\) lifts to a morphism \(g: Y' \to X\) such that \(g|_{Y} = f\).
\end{proposition}

\begin{proof}
    First we note that $ Y' $ is also affine~\cite[\uppercase\expandafter{\romannumeral3}, Ex. 3.1]{Har77}, so we may transfer the problem to algebraic setting. Let \(f:A\to B\) be the ring homomorphism over \(k\) corresponding to \(f:Y\to X\) where \(B = B^{\prime} / I\) with \(I^{n} = 0\) for some \(n\). We want to find a homomorphism \(g:A\to B^{\prime}\) lifting \(f\).

    Consider the sequence \(B^{\prime} = B^{\prime} / I^{n}\to B^{\prime} / I^{n-1}\to \dots \to B^{\prime} / I^{2}\to B^{\prime} / I\), it will suffice to lift one step at a time. So we may assume that \(I^{2} = 0\).

    Since \(X\) is of finite type over \(k\), we may write \(A = k[x_{1},\ldots ,x_{n}] / J\). Consider any lifting \(h:k[x_{1},\ldots ,x_{n}]\to B^{\prime}\) of \(k[x_{1},\ldots ,x_{n}]\to A\to B\). Then \(h\) induces a map \(J\to I\), and thus induces \(\overline{h}:J / J^{2}\to I\) since \(I^{2} = 0\).

    Embed \(X\hookrightarrow \mathbb{A}_{k}^{n}\) by \(P = k[x_{1},\ldots ,x_{n}]\to A\). Since \(X\) is smooth, we obtain an exact sequence\[
        0\longrightarrow J / J^{2}\longrightarrow \varOmega_{P / k}^{1}\otimes_{P}A\longrightarrow \varOmega_{A / k}^{1}\longrightarrow 0.
    \]Notice that those modules correspond to locally free sheaves on \(X\), and thus are projective \(A\)-modules. Via the maps \(h\), \(f\), we get a \(P\)-module structure on \(B'\), and \(A\)-module structures on \(B\), \(I\). Applying \(\mathrm{Hom}_{A}(-, I)\) will get another exact sequence\[
        0\longrightarrow \mathrm{Hom}_{A}(\varOmega_{A / k}^{1},I)\longrightarrow \mathrm{Hom}_{P}(\varOmega_{P / k}^{1},I)\longrightarrow \mathrm{Hom}_{A}(J / J^{2},I)\longrightarrow 0.
    \]Let \(\theta \in \mathrm{Hom}_{P}(\varOmega_{P / k}^{1},I) = \mathrm{Der}_{k}(P,I)\) be an element whose image is \(\overline{h}\in \mathrm{Hom}_{A}(J / J^{2},I)\). We claim that \(h^{\prime}:= h - \theta :P\to B^{\prime}\) is a ring homomorphism lifting \(f\) with \(h^{\prime}(J) = 0\). The first assertion follows by the following lemma.

\begin{lemma}\label{A14}
    Let \(B^{\prime}\to B\) be a surjective homomorphism of \(k\)-algebras with kernel \(I\) of square zero. Let \(R\to B\) be a homomorphism of \(k\)-algebras.
    \begin{enumerate}[(1)]
        \item If \(f,g: R \to B'\) are two liftings of the map \(R \to B\) to \(B'\), then \(\theta = g - f\) is a \(k\)-derivation of \(R\) to \(I\).
        \item Conversely, if \(f: R \to B'\) is one lifting, and \(\theta: R \to I\) is a derivation, then \(g = f + \theta\) is another homomorphism of \(R\) to \(B'\) lifting the given map \(R \to B\).
    \end{enumerate}
    In other words, if it is nonempty, the set of liftings \(R \to B\) to \(k\)-algebra homomorphisms of \(R \to B^{\prime}\) is a principal homogeneous space under the action of \(\mathrm{Der}_{k}(R,I) = \mathrm{Hom}_{R}(\varOmega_{R / k},I)\).
\end{lemma}

\begin{proof}
    \textit{(1)} Let $ f,g:R\rightarrow B' $ and let $ \theta=g-f $. As a $ k $-linear map, $ \theta $ followed by the projection $ B'\rightarrow B $ is zero, so $ \theta $ sends $ R $ to $ I $. Let $ x,y\in R $. Then\begin{align*}
        \theta(xy)&=g(xy)-f(xy)\\
        &=g(x)g(y)-f(x)f(y)\\
        &=g(x)(g(y)-f(y))+f(y)(g(x)-f(x))\\
        &=g(x)\theta(y)+f(y)\theta(x)\\
        &=x\theta(y)+y\theta(x),
    \end{align*}the last step being because $ g(x) $ and $ f(y) $ act in $ I $ just like $ x,y $. Thus $ \theta $ is a $ k $-derivation of $ R $ to $ I $.

    \textit{(2)} Conversely, given $ f $ and $ \theta $ as above, let $ g=f+\theta $. Then\begin{align*}
        g(xy)&=f(xy)+\theta(xy)\\
        &=f(x)f(y)+x\theta(y)+y\theta(x)\\
        &=(f(x)+\theta(x))(f(y)+\theta(y))\\
        &=g(x)g(y),
    \end{align*}where we note that $ \theta(x)\theta(y)=0 $, since $ I^2=0 $. Thus $ g $ is a homomorphism of $ R\rightarrow B' $ lifting $ R $ to $ B $.
\end{proof}

    Since \(h^{\prime}(J)\subseteq I\), we have \(h^{\prime}(J^{2})\subseteq I^{2} = 0\). Then \(h^{\prime}\) descends to \(P / J^{2}\to B^{\prime}\), and thus \(h^{\prime}|_{J / J^{2}} = h - h = 0\) by the choice of \(\theta\). Hence, \(h^{\prime}\) descends to the desired homomorphism \(g:A\to B^{\prime}\) lifting \(f\).
\end{proof}

Next, we investigate the relation between nonsingularity and the \(T^{i}\) functors.

\begin{theorem}\label{A17}
    Let \(X = \mathrm{Spec}\,  B\) be an affine scheme over an algebraically closed field \(k\). Then \(X\) is nonsingular if and only if \(T^{1}(B / k,M) = 0\) for all \(M\in \mathsf{Mod}_{B}\). Furthermore, if \(X\) is nonsingular, then \(T^{2}(B / k,M) = 0\) for all \(M\).
\end{theorem}

\begin{proof}
    Write \(B\) as a quotient of a polynomial ring \(A = k[x_{1},\ldots ,x_{n}]\) over \(k\). Then \(X\) is nonsingular if and only if the conormal sequence\begin{equation}\label{eq12}
        0\longrightarrow I / I^{2}\longrightarrow \varOmega_{A / k}\otimes_{A}B\longrightarrow \varOmega_{B / k}\longrightarrow 0
    \end{equation}is exact and \(\varOmega_{B / k}\) is locally free, i.e., a projective \(B\)-module. Since \(\varOmega_{A / k}\) is a free \(A\)-module, we conclude that \(X\) is nonsingular if and only if (\ref{eq12}) is split exact. From the four-term exact sequence in (\ref{A9})\[
        0\longrightarrow T^{0}(B / k,M)\longrightarrow \mathrm{Hom}(\varOmega_{A / k},M)\longrightarrow \mathrm{Hom}(I / I^{2},M)\longrightarrow T^{1}(B / k,M)\longrightarrow 0,
    \]\(T^{1}(B / k,M) = 0\) for all \(M\in \mathsf{Mod}_{B}\) if and only if the map\[
        \mathrm{Hom}(\varOmega_{A / k},M)\longrightarrow \mathrm{Hom}(I / I^{2},M)
    \]is surjective for all \(M\), which is equivalent to split exactness by considering \(M = I / I^{2}\).

    Again, by (\ref{A9}) we have \(T^{2}(B / k,M) = T^{2}(B / A,M)\). Since \(X\) is smooth, the localizing ideal \(I_{x}\) at \(x\in X\) is generated by a regular sequence in the regular local ring \(A_{x}\). By (\ref{A8}), \(T^{2}(B_{x} / A_{x},M) = 0\) for all \(M\in \mathsf{Mod}_{B_{x}}\). By (\ref{A11}), we conclude that \(T^{2}(B / A,M) = 0\) for all \(M\in \mathsf{Mod}_{B}\).
\end{proof}

\begin{theorem}\label{A18}
    A morphism of finite type \(f:X\to Y\) of Noetherian schemes is smooth if and only if it is flat, and \(\mathcal{T}^{1}(X / Y,\mathcal{F}) = 0\) for all coherent sheaves \(\mathcal{F}\in \mathsf{Coh}(X)\). Furthermore, if \(f\) is smooth, then \(\mathcal{T}^{2}(X / Y,\mathcal{F}) = 0\) for all \(\mathcal{F}\).
\end{theorem}

\begin{proof}
    The problem is local, so we may assume that \(X = \mathrm{Spec}\,  B\) and \(Y = \mathrm{Spec}\,  A\) are affine and \(f\) is given by a ring homomorphism \(A\to B\).

    First suppose \(B\) is flat over \(A\) and \(T^{1}(B / A,M) = 0\) for all \(M\in \mathsf{Mod}_{B}\). Let \(y\in Y\) correspond to \(\mathfrak{p}\in \mathrm{Spec}\,  A\) and let \(k = k(y)\) be its residue field. Let \(A^{\prime} = A / \mathfrak{p}\) and \(B^{\prime} = B\otimes_{A}A^{\prime} = B / \mathfrak{p}B\); then for any \(M\in \mathsf{Mod}_{B^{\prime}}\) we obtain \(T^{1}(B^{\prime} / A^{\prime},M) = T^{1}(B / A,M) = 0\) by (\ref{A12}). Write \(B^{\prime} = A^{\prime}[x_{1},\ldots ,x_{n}] / I\); then \(T^{1}(B^{\prime} / A^{\prime},I / I^{2}) = 0\).

    Now consider the flat base extension from \(A^{\prime}\) to \(\overline{k} = \overline{\mathrm{Frac}(A^{\prime})}\). By (\ref{A12}), we have\[
        T^{1}(B^{\prime}\otimes_{A^{\prime}}\overline{k} /\overline{k},(I / I^{2})\otimes \overline{k}) = 0.
    \]Since \(\overline{k}\) is flat over \(A^{\prime}\), we have \(B^{\prime}\otimes \overline{k} = \overline{k}[x_{1},\ldots ,x_{n}] / \overline{I}\) and \((I / I^{2})\otimes \overline{k} = \overline{I} / \overline{I}^{2}\), where \(\overline{I} = I\otimes \overline{k}\). Then the proof of (\ref{A17}) shows that \(\mathrm{Spec}\,  B^{\prime}\otimes \overline{k}\) is nonsingular over \(\overline{k}\). Hence, the geometric fibers of \(f\) are nonsingular, and thus \(f\) is smooth.

    Conversely, suppose that \(B\) is smooth over \(A\). We claim that \(T^{1}(B / A,B / \mathfrak{m}) = 0\) for all maximal ideals \(\mathfrak{m}\) of \(B\). Indeed, let \(\mathfrak{m}\) correspond to the point \(x\in \mathrm{Spec}\,  B\), let \(f(x) = y\), and let \(k = k(y)\). Apply base change on \(B / \mathfrak{m}\in \mathsf{Mod}_{B\otimes_{A}k}\), we obtain\[
        T^{1}(B / A,B / \mathfrak{m})\otimes_{k}\overline{k} = T^{1}(B\otimes_{A}k / k,B / \mathfrak{m})\otimes_{k}\overline{k} = T^{1}(B \otimes_{A}\overline{k} /\overline{k},(B / \mathfrak{m})\otimes_{k}\overline{k}) = 0
    \]since the geometric fibers are nonsingular. Since \(k\to \overline{k}\) is faithfully flat, we conclude that \(T^{1}(B / A,B / \mathfrak{m}) = 0\).
\end{proof}

\subsection{Deformations of schemes and obstruction theory}

Let $ k $ be a fixed algebraically closed ground field, and let $ \mathsf{C} $ be the category of local artinian $ k $-algebras with residue field $ k $. Let $ X_0 $ be a scheme over $ k $ and take a $ C\in\mathsf{C} $. 

A \textbf{deformation} of $ X_0 $ over $ C $ is a pair $ (X,i) $, where $ X $ is a scheme flat over $ C $, and $ i:X_0\rightarrow X $ is a closed immersion such that the induced map $ i\otimes k:X_0\rightarrow X\otimes_C k $ is an isomorphism.

Let $ X_0\hookrightarrow X $ be a deformation of $ X_0 $ over $ C\in \mathsf{C} $. Suppose $ C' $ is another artinian $ k $-algebra with residue field $ k $, equipped with a surjective map $ C'\rightarrow C $ with the kernel $ J $ annihilated by $ \mathfrak{m}_{C'} $.

An \textbf{extension} of $ X $ over $ C' $ is a deformation $ X' $ of $ X_0 $ over $ C' $, together with a closed immersion $ X\hookrightarrow X' $ inducing an isomorphism $ X\stackrel{\sim}{\longrightarrow}X'\times_{C'}C $. Two such extension $ X' $ and $ X'' $ are \textbf{equivalent} if there is an isomorphism of deformations $ X'\stackrel{\sim}{\longrightarrow}X'' $ compatible with the respective closed immersion of $ X $ into $ X' $ and $ X'' $. Let $ \mathrm{Def}(X/C,C') $ be the set of equivalence classes of such extensions of $ X' $ over $ C' $. In particular, define two deformations of $ X_0 $ over $ C $ are equivalent if they are in the same class in $ \mathrm{Def}(X_0/k,C) $.

% \begin{remark}
%     There is a surjective map\[
%         \mathrm{Def}(X/C,C')\longrightarrow \mathrm{Def}(C'\rightarrow C)^{-1}(X),
%     \]which is not injective in general, since there may exist an automorphism of $ X' $ that restricts to the identity on $ X_0 $ but does not extend to an automorphism of $ X $.
% \end{remark}

\begin{theorem}\label{B2}
    In the situation described as above, with the additional structure as follows: let $ Y_0\subseteq X_0 $ and $ Y\subseteq X $ be closed subschemes such that $ Y_0=Y\times_C k $ under the identification $ X_0\stackrel{\sim}{\longrightarrow}X\times_C k $. Fixed an extension $ X' $ of $ X $ over $ C' $. An \textbf{(embedding) extension} of $ Y $ over $ C' $ in $ X' $ is a closed subscheme $ Y'\subseteq X' $, flat over $ C' $, such that $ Y'\times_{C'}C=Y $. Then the set of such extension of $ Y $ over $ C' $ in $ X' $ form a $ H^0(Y_0,\mathcal{N}_{Y_0/X_0}\otimes_k J) $-pseudotorsor.
\end{theorem}

\begin{proof}
    First we consider the affine case $ X=\mathrm{Spec}\, B $, $ X'=\mathrm{Spec}\, B' $, $ Y=\mathrm{Spec}\, A $, and $ A=B/I $. Suppose that $ Y'=\mathrm{Spec}\, A' $ is given by $ A'=B'/I' $. Tensoring the short exact sequence\[
        0\longrightarrow I'\longrightarrow B'\longrightarrow A'\longrightarrow 0
    \]with the sequence $ 0\rightarrow J\rightarrow C'\rightarrow C\rightarrow 0 $, we will get a diagram% https://quiver.taskfirst.cn/#q=WzAsMjEsWzEsMCwiMCJdLFsyLDAsIjAiXSxbMywwLCIwIl0sWzQsMSwiMCJdLFs0LDIsIjAiXSxbNCwzLCIwIl0sWzEsNCwiMCJdLFsyLDQsIjAiXSxbMyw0LCIwIl0sWzAsMSwiMCJdLFswLDIsIjAiXSxbMCwzLCIwIl0sWzEsMSwiSlxcb3RpbWVzX0NJIl0sWzIsMSwiSSciXSxbMywxLCJJIl0sWzEsMiwiSlxcb3RpbWVzX0NCIl0sWzIsMiwiQiciXSxbMywyLCJCIl0sWzEsMywiSlxcb3RpbWVzX0NBIl0sWzIsMywiQSciXSxbMywzLCJBIl0sWzksMTJdLFsxMiwxM10sWzEzLDE0XSxbMTQsM10sWzEwLDE1XSxbMTUsMTZdLFsxNiwxN10sWzE3LDRdLFswLDEyXSxbMSwxM10sWzIsMTRdLFsxMiwxNV0sWzEzLDE2XSxbMTQsMTddLFsxMSwxOF0sWzE4LDE5XSxbMTksMjBdLFsyMCw1XSxbMTUsMThdLFsxOCw2XSxbMTksN10sWzE2LDE5XSxbMTcsMjBdLFsyMCw4XV0=
    \[\begin{tikzcd}
        & 0 & 0 & 0 & \\
        0 & {J\otimes_CI} & {I'} & I & 0 \\
        0 & {J\otimes_CB} & {B'} & B & 0 \\
        0 & {J\otimes_CA} & {A'} & A & 0 \\
        & 0 & 0 & 0
        \arrow[from=1-2, to=2-2]
        \arrow[from=1-3, to=2-3]
        \arrow[from=1-4, to=2-4]
        \arrow[from=2-1, to=2-2]
        \arrow[from=2-2, to=2-3]
        \arrow[from=2-2, to=3-2]
        \arrow[from=2-3, to=2-4]
        \arrow[from=2-3, to=3-3]
        \arrow[from=2-4, to=2-5]
        \arrow[from=2-4, to=3-4]
        \arrow[from=3-1, to=3-2]
        \arrow[from=3-2, to=3-3]
        \arrow[from=3-2, to=4-2]
        \arrow[from=3-3, to=3-4]
        \arrow[from=3-3, to=4-3]
        \arrow[from=3-4, to=3-5]
        \arrow[from=3-4, to=4-4]
        \arrow[from=4-1, to=4-2]
        \arrow[from=4-2, to=4-3]
        \arrow[from=4-2, to=5-2]
        \arrow[from=4-3, to=4-4]
        \arrow[from=4-3, to=5-3]
        \arrow[from=4-4, to=4-5]
        \arrow[from=4-4, to=5-4]
    \end{tikzcd}\]By (\ref{wto}), the exactness of the bottom two rows is equivalent to the flatness of $ B' $ and $ A' $ over $ C $, respectively. The exactness of the first column follows from the flatness of $ A $ over $ C $. Now we seek to classify the possible $ A' $ (resp. $ I' $) full in the diagram.

    Suppose that $ I' $ and $ I'' $ are two choices of $ I' $ to fill in the diagram. Given $ x\in I $, lift it to $ x'\in I' $ and $ x''\in I'' $. Then $ x''-x'\in B' $ maps to $ 0 $ in $ B $, and thus $ x''-x'\in J\otimes_C B $. Denote its image in $ J\otimes_C A $ by $ \varphi(x) $. Notice that the choice of $ x' $ and $ x'' $ may differ by elements in $ J\otimes_C I $, which maps to $ 0 $ in $ J\otimes_C A $. So $ \varphi $ is well-defined additive map. In fact, it is an $ B $-linear homomorphism $ \varphi\in \mathrm{Hom}_B(I,J\otimes_C A) $.
    
    Conversely, given $ I' $ and $ \varphi\in \mathrm{Hom}_B(I,J\otimes_C A) $, we can define another ideal $ I'' $ solving the extension problem as follows. $ I'' $ is the set of $ x''\in B' $ whose image in $ B $ is $ x\in I $, and such that for any lifting $ x'\in I' $ of $ x $, the image $ x''-x $ in $ J\otimes_C A $ is equal to $ \varphi(x) $.

    If $ I' $, $ I'' $, $ I''' $ are three choice of solution, and if $ \varphi_1 $ is defined by $ I' $, $ I'' $ as above, $ \varphi_2 $ defined by $ I'' $, $ I''' $, and $ \varphi_3 $ defined by $ I' $, $ I''' $, then $ \varphi_3=\varphi_1+\varphi_2 $. So we conclude that\begin{align*}
        \mathrm{Hom}_B(I,J\otimes_C A)\times\{\text{deformation of}\ Y\ \text{over}\ C'\ \text{in}\ X'\}& \longrightarrow\{\text{deformation of}\ Y\ \text{over}\ C'\ \text{in}\ X'\}\\
        (I',\varphi)& \longmapsto I''
    \end{align*}is an proper transitive group action. Since we can't guarantee the existence, it only a pseudotorsor. In this case,\[
        \mathrm{Hom}_B(I,J\otimes_C A)=\mathrm{Hom}_B(I,J\otimes_k A_0).
    \]Since $ J $ is a $ k $-vector space, this term becomes to\[
        \mathrm{Hom}_{B_0}(I_0,J\otimes_k A_0)=\mathrm{Hom}_{A_0}(I_0\otimes_{B_0}A_0,A_0\otimes_k J)=H^0(Y_0,\mathcal{N}_{Y_0/X_0}\otimes_k J).
    \]In general case, the group action $ H^0(Y_0,\mathcal{N}_{Y_0/X_0}\otimes_k J) $ on the set of $ Y' $ can be defined locally and functorial. So we may glue it to the global action.
\end{proof}

Now back to deformations of schemes (rather than closed subschemes). Consider the affine scheme case first. Given a finite generated $ k $-algebra $ B_0 $, and a deformation $ B $ of $ B_0 $ over $ C $. We want to describe the set $ \mathrm{Def}(B/C,C'):=\mathrm{Def}(\mathrm{Spec}\, B/C,C') $.

\begin{theorem}\label{B3}
    In the situation as beginning:\begin{enumerate}[(1)]
        \item There is an element $ \delta\in T^2(B_0/k,B_0\otimes J) $, called the \textbf{obstruction}, with the property that $ \delta=0 $ iff an extension $ B' $ of $ B $ exists.
        \item If extension exists, then $ \mathrm{Def}(B/C,C') $ form a $ T^1(B_0/k,B_0\otimes J) $-torsor.
    \end{enumerate}
\end{theorem}

\begin{proof}
    Take a presentation $ R=C[x_1,\dots,x_n]\rightarrow B $ with kernel $ I=(f_1,\dots,f_r) $. Define a morphism $ F:=R^r\rightarrow I $ by $ e_i\mapsto f_i $ and denoted the kernel by $ Q $. Lift each $ f_i $ to an element $ f'_i\in R':=C'[x_1,\dots,x_n] $, and define $ I'=(f'_1,\dots,f'_r) $, $ B'=R'/I' $, $ F'=R'^r $. Now consider\begin{align*}
        0\longrightarrow I'\longrightarrow R'\longrightarrow B'\longrightarrow 0 \qquad 0\longrightarrow Q'\longrightarrow & F'\longrightarrow I'\longrightarrow 0.\\
        & e'_i\longmapsto f'_i
    \end{align*}Tensoring $ C $ over $ C' $, we have\[
        B'\otimes_{C'}C\cong\mathrm{coker}(I'\otimes_{C'}C\rightarrow R'\otimes_{C'}C)=B.
    \]Tensoring with $ 0\rightarrow J\rightarrow C'\rightarrow C\rightarrow 0 $ over $ C' $, we get a diagram% https://quiver.taskfirst.cn/#q=WzAsMjEsWzAsMiwiMCJdLFsxLDIsIkZcXG90aW1lc19DIEoiXSxbMiwwLCIwIl0sWzMsMCwiMCJdLFsyLDEsIlEnIl0sWzMsMSwiUSJdLFsyLDIsIkYnIl0sWzMsMiwiRiJdLFs0LDIsIjAiXSxbNCwzLCIwIl0sWzQsNCwiMCJdLFszLDUsIjAiXSxbMiw1LCIwIl0sWzEsNSwiMCJdLFswLDMsIjAiXSxbMSwzLCJSXzBcXG90aW1lc19rIEoiXSxbMSw0LCJCXzBcXG90aW1lc19rIEoiXSxbMiw0LCJCJyJdLFszLDQsIkIiXSxbMywzLCJSIl0sWzIsMywiUiciXSxbMiw0XSxbMyw1XSxbNCw2XSxbNSw3XSxbNyw4XSxbNiw3XSxbMSw2XSxbMCwxXSxbMTQsMTVdLFsxNSwyMF0sWzE2LDE3XSxbMTcsMThdLFsxOCwxMF0sWzE5LDldLFsyMCwxOV0sWzYsMjAsImYnIl0sWzcsMTksImYiXSxbMSwxNV0sWzE1LDE2XSxbMTYsMTNdLFsyMCwxN10sWzE3LDEyXSxbMTksMThdLFsxOCwxMV1d
    \[\begin{tikzcd}
        && 0 & 0 & \\
        && {Q'} & Q \\
        0 & {F\otimes_C J} & {F'} & F & 0 \\
        0 & {R_0\otimes_k J} & {R'} & R & 0 \\
        & {B_0\otimes_k J} & {B'} & B & 0 \\
        & 0 & 0 & 0
        \arrow[from=1-3, to=2-3]
        \arrow[from=2-3, to=2-4]
        \arrow[from=1-4, to=2-4]
        \arrow[from=2-3, to=3-3]
        \arrow[from=2-4, to=3-4]
        \arrow[from=3-1, to=3-2]
        \arrow[from=3-2, to=3-3]
        \arrow[from=3-2, to=4-2]
        \arrow[from=3-3, to=3-4]
        \arrow["{f'}", from=3-3, to=4-3]
        \arrow[from=3-4, to=3-5]
        \arrow["f", from=3-4, to=4-4]
        \arrow[from=4-1, to=4-2]
        \arrow[from=4-2, to=4-3]
        \arrow[from=4-2, to=5-2]
        \arrow[from=4-3, to=4-4]
        \arrow[from=4-3, to=5-3]
        \arrow[from=4-4, to=4-5]
        \arrow[from=4-4, to=5-4]
        \arrow[from=5-2, to=5-3]
        \arrow[from=5-2, to=6-2]
        \arrow[from=5-3, to=5-4]
        \arrow[from=5-3, to=6-3]
        \arrow[from=5-4, to=5-5]
        \arrow[from=5-4, to=6-4]
    \end{tikzcd}\]where we use $ -\otimes_{C'}J=-\otimes_{C'}k\otimes_{k}J $ and $ R_0=k[x_1,\dots,x_n] $. By snake lemma, there is a map $ \delta_0:Q\rightarrow B_0\otimes_k J $ depending on the lifting $ f' $ of $ f $. By (\ref{wto}), $ B' $ is flat over $ C' $ iff $ \delta_0=0 $.

    Recall that in the definition of cotangent complex, any element in $ F_0 $ is of the form $ f_j e_i-f_i e_j\in F $, and lifts to $ f_j' e_i'-f_i'e_j'\in Q' $. So the map $ \delta_0 $ factors through $ \delta_1:Q/F_0\rightarrow B_0\otimes_k J $. By definition of $ T^i $ functor and base change (\ref{A12}), we get\begin{align*}
        \mathrm{coker}(\mathrm{Hom}(F/IF,B_0\otimes_k J)&\longrightarrow \mathrm{Hom}(Q/F_0,B_0\times_k J))=T^2(B/C,B_0\otimes_k J)=T^2(B_0/k,B_0\otimes_k J)\\
        \bar{\delta}_1 &\longmapsto\delta.
    \end{align*}Here $ \delta\in T^2(B_0/k,B_0\otimes_k J) $.

    We claim that $ \delta $ is independent of all choices made. If we make a different choice of lifting $ f_i'' $ of the $ f_i' $, then $ f_i'-f_i'' $ define a map from $ F' $ to $ R_0\otimes_k J $, and hence from $ F/IF\rightarrow B_0\otimes_k J $, and these go to zero in $ T^2 $. Suppose we choose a different polynomial ring $ R^*\rightarrow B $. As in the proof of (\ref{A1}), we may reduce to the case $ R^*=R[y_1,\dots,y_s] $, and $ y_i\mapsto 0\in B $. Then the $ \delta_0^*(y_i)=0 $ implies $ \delta=\delta^* $ in $ T^2(B_0/k,B_0\otimes_k J) $.

    If the extension $ B'/C' $ exists, take a polynomial ring $ R' $ over $ C' $ maps surjectively to $ B' $. Consider\begin{equation}\label{sgh}
        0\longrightarrow I'\longrightarrow R'\longrightarrow B'\longrightarrow 0 \qquad 0\longrightarrow Q'\longrightarrow F'\longrightarrow I'\longrightarrow 0.
    \end{equation}Since $ B' $ and $ R' $ are flat over $ C' $, we also have $ I' $ is flat over $ C' $. Let $ R=R'\otimes_{C'}C $, $ I=I'\otimes_{C'}C $, and $ Q=Q'\otimes_{C'}C $. Then (\ref{sgh})$ \otimes_{C'}C $ will give a presentation to calculate cotangent complex of $ B $ over $ C $. Since $ Q'\rightarrow Q $ is surjective, it follows that $ \delta_0=0 $, and thus $ \delta=0 $.

    Conversely, suppose that $ \delta=0 $. We claim that for a suitable choose of lifting $ f_i' $, the induced map $ \delta_0:Q\rightarrow B_0\otimes_k J $ is zero. By assumption, the map $ \delta_1\in\mathrm{Hom}(Q/F_0,B_0\otimes_k J) $ lifts to a map $ \gamma:F/IF\rightarrow B_0\otimes_k J $, and defines a map $ F\rightarrow B_0\otimes_k J $. Since $ F $ is free, it lifts to a map $ F\rightarrow R_0\otimes_k J $. Denote the image of $ e_i $ by $ g_i\in R_0\otimes_k J $. Now take $ f_i''=f'_i-g_i $, then the new $ \delta_0 $ is zero, and thus the new $ B' $ is flat over $ C' $. This complete the proof of \textit{(1)}.

    Suppose one such extension $ B_1' $ exists. Take a presentation $ 0\rightarrow I_0\rightarrow R_0\rightarrow B_0\rightarrow 0 $ as usual. Lift the surjection $ R_0\rightarrow B_0 $ to a map $ R\rightarrow B $, and to a map $ R'\rightarrow B_1' $, which is compatible with $ B_1'\rightarrow B\rightarrow B_0 $. By Nakayama's lemma, $ R\rightarrow B $ and $ R'\rightarrow B_1' $ are surjective.

    For any other extension $ B_2' $ of $ B $, the map $ R\rightarrow B $ lifts to a map $ R'\rightarrow B_2' $. Thus, every abstract deformation is also an embedded deformation $ X\hookrightarrow\mathbb{A}^n_k $. By (\ref{B2}), the embedded deformation is a $ \mathrm{Hom}_B(I/I^2,B\otimes_C J) $-torsor. Suppose that we have two embedding deformations $ B_2' $ and $ B_3' $ are equivalent as abstract extensions of $ B $. Choose an isomorphism $ B_2'\cong B_3' $ as extension of $ B $ over $ C' $. Then we obtain two maps $ R'\rightarrow B_2' $, which can be regarded as the lifting of $ R'\rightarrow R\rightarrow B $. By (\ref{A14}), the differ of this two map is a $ C' $-derivation of $ R' $ to $ B_2'\otimes_{C'}J\cong B\otimes_C J $, which can be regarded as an element of\[
        \mathrm{Hom}_{R'}(\varOmega_{R'/C'},B\otimes_C J)=\mathrm{Hom}_R(\varOmega_{R'/C'}\otimes_{C'}C,B\otimes_C J)=\mathrm{Hom}_R(\varOmega_{R/C},B\otimes_C J).
    \]Using the exact sequence (\ref{A9})\[
        \mathrm{Hom}_R(\varOmega_{R/C},B\otimes_C J)\longrightarrow \mathrm{Hom}(I/I^2,B\otimes_C J)\longrightarrow T^1(B/C,B\otimes_C J)\longrightarrow 0,
    \]we observe that the ambiguity of embedding is exactly resolved by the image of the derivations. We conclude that $ \mathrm{Def}(X/C,C') $ is a torsor under $ T^1(B/C,B_0\otimes_k J)=T^1(B_0/k,B_0\otimes_k J) $ by (\ref{A12}).
\end{proof}

\begin{remark}\label{B4}
    Given an extension $ B' $ of $ B $ over $ C' $, the automorphism group of $ B' $ as an extension of $ B $ over $ C $ is naturally isomorphic to the group $ T^0(B_0/k,J\otimes_k B_0) $. Indeed, take $ k=C' $ and $ R=B' $ in (\ref{A14}), we have isomorphism\begin{align*}
        \mathrm{Der}_{C'}(B',J\otimes_{C'}B')&\stackrel{\sim}{\longrightarrow}\mathrm{Aut}(B'/B)=\mathrm{Der}_{C'}(B',B)\\
        \theta&\longmapsto \mathrm{id}_{B'}+\theta
    \end{align*}where $ \theta_1+\theta_2\mapsto (\mathrm{id}_{B'}+\theta_1)(\mathrm{id}_{B'}+\theta_2)=\mathrm{id}_{B'}+\theta_1+\theta_2 $ since $ J\otimes_{C'}B' $ is square zero ideal. Moreover, we have\begin{align*}
        \mathrm{Der}_{C'}(B',J\otimes_{C'}B')&=\mathrm{Hom}_{B'}(\varOmega_{B'/C'},J\otimes_{C'}B')=\mathrm{Hom}_{B_0}(\varOmega_{B'/C'}\otimes_{B'}B_0,J\otimes_k B_0)\\
        &=\mathrm{Hom}_{B_0}(\varOmega_{B_0/k},J\otimes_k B_0)=T^0(B_0/k,B_0\otimes_k J).
    \end{align*}
\end{remark}

For general case, we recall the notation $ \mathcal{T}^i_{X_0}=\mathcal{T}^i(X_0/k,\mathcal{O}_{X_0}) $.

\begin{lemma}[local criterion of flatness]\label{wto}
    Let $ A'\rightarrow A $ be a surjective homomorphism of noetherian rings whose kernel $ J $ is square zero. Then an $ A' $-module $ M' $ is flat over $ A' $ iff\begin{enumerate}[(1)]
        \item $ M=M'\otimes_{A'}A $ is flat over $ A $
        \item the natural map $ M\otimes_A J\rightarrow M' $ is injective.
    \end{enumerate}
\end{lemma}

\begin{proof}
    Since $ J^2=0 $, we may regard $ J $ as $ A $-module and identify $ M'\otimes_{A'}J $ with $ M\otimes_A J $.

    If $ M' $ is flat over $ A' $, then \textit{(1)} follows by base extension and \textit{(2)} follows by tensoring $ M' $ with the exact sequence\[
        0\longrightarrow J\longrightarrow A'\longrightarrow A\longrightarrow 0.
    \]

    Conversely, suppose that $ M' $ satisfies \textit{(1)} and \textit{(2)}. It suffices to prove $ \mathrm{Tor}^{A'}_1(M',A'/\mathfrak{p}')=0 $ for every $ \mathfrak{p}'\subseteq A' $. Since $ J $ is nilpotent, it is contained in $ \mathfrak{p}' $. Let $ \mathfrak{p}=\mathfrak{p}'/J\in \mathrm{Spec}\, A $, then we have the following diagram% https://quiver.taskfirst.cn/#q=WzAsMTYsWzAsMSwiMCJdLFswLDIsIjAiXSxbNCwxLCIwIl0sWzQsMiwiMCJdLFsyLDQsIjAiXSxbMyw0LCIwIl0sWzIsMCwiMCJdLFszLDAsIjAiXSxbMSwxLCJKIl0sWzEsMiwiSiJdLFsyLDEsIlxcbWF0aGZyYWt7cH0nIl0sWzMsMSwiXFxtYXRoZnJha3twfSJdLFsyLDIsIkEnIl0sWzMsMiwiQSJdLFsyLDMsIkEnL1xcbWF0aGZyYWt7cH0nIl0sWzMsMywiQS9cXG1hdGhmcmFre3B9Il0sWzAsOF0sWzgsOSwiIiwxLHsibGV2ZWwiOjIsInN0eWxlIjp7ImhlYWQiOnsibmFtZSI6Im5vbmUifX19XSxbMSw5XSxbOSwxMl0sWzEyLDEzXSxbMTMsM10sWzEyLDE0XSxbMTMsMTVdLFsxMCwxMl0sWzYsMTBdLFs3LDExXSxbMTEsMTNdLFsxNSw1XSxbMTQsNF0sWzgsMTBdLFsxMCwxMV0sWzExLDJdXQ==
    \[\begin{tikzcd}
        && 0 & 0 & \\
        0 & J & {\mathfrak{p}'} & {\mathfrak{p}} & 0 \\
        0 & J & {A'} & A & 0 \\
        && {A'/\mathfrak{p}'} & {A/\mathfrak{p}} \\
        && 0 & 0
        \arrow[from=1-3, to=2-3]
        \arrow[from=1-4, to=2-4]
        \arrow[from=2-1, to=2-2]
        \arrow[from=2-2, to=2-3]
        \arrow[equals, from=2-2, to=3-2]
        \arrow[from=2-3, to=2-4]
        \arrow[from=2-3, to=3-3]
        \arrow[from=2-4, to=2-5]
        \arrow[from=2-4, to=3-4]
        \arrow[from=3-1, to=3-2]
        \arrow[from=3-2, to=3-3]
        \arrow[from=3-3, to=3-4]
        \arrow[from=3-3, to=4-3]
        \arrow[from=3-4, to=3-5]
        \arrow[from=3-4, to=4-4]
        \arrow[from=4-3, to=5-3]
        \arrow[from=4-4, to=5-4]
    \end{tikzcd}\]Tensoring with $ M' $, we obtain% https://quiver.taskfirst.cn/#q=WzAsMTYsWzEsMSwiMCJdLFsyLDEsIjAiXSxbMywyLCIwIl0sWzMsMywiMCJdLFsxLDUsIjAiXSxbMiw1LCIwIl0sWzEsMiwiXFxidWxsZXQiXSxbMiwyLCJcXGJ1bGxldCJdLFsyLDMsIk0iXSxbMiw0LCJNIl0sWzEsNCwiTSciXSxbMSwzLCJNJyJdLFswLDMsIk1cXG90aW1lc19BIEoiXSxbMCwyLCJNXFxvdGltZXNfQSBKIl0sWzEsMCwiMCJdLFsyLDAsIjAiXSxbMTMsMTIsIiIsMSx7ImxldmVsIjoyLCJzdHlsZSI6eyJoZWFkIjp7Im5hbWUiOiJub25lIn19fV0sWzEyLDExXSxbMTEsOF0sWzgsM10sWzgsOV0sWzksNV0sWzExLDEwXSxbMTAsNF0sWzE0LDBdLFsxNSwxXSxbMCw2XSxbMSw3XSxbMCwxXSxbNiw3XSxbNywyXSxbMTMsNl0sWzYsMTFdLFs3LDhdXQ==
    \[\begin{tikzcd}
        & 0 & 0 & \\
        & \mathrm{Tor}^{A'}_1(M',A'/\mathfrak{p}') & \mathrm{Tor}^{A}_1(M,A/\mathfrak{p}) \\
        {M\otimes_A J} & M'\otimes_{A'}\mathfrak{p}' & M\otimes_A\mathfrak{p} & 0 \\
        {M\otimes_A J} & {M'} & M & 0 \\
        & {M'\otimes_{A'}A'/\mathfrak{p}'} & M\otimes_{A}A/\mathfrak{p} \\
        & 0 & 0
        \arrow[from=1-2, to=2-2]
        \arrow[from=1-3, to=2-3]
        \arrow[from=2-2, to=2-3]
        \arrow[from=2-2, to=3-2]
        \arrow[from=2-3, to=3-3]
        \arrow[from=3-1, to=3-2]
        \arrow[equals, from=3-1, to=4-1]
        \arrow[from=3-2, to=3-3]
        \arrow[from=3-2, to=4-2]
        \arrow[from=3-3, to=3-4]
        \arrow[from=3-3, to=4-3]
        \arrow[from=4-1, to=4-2]
        \arrow[from=4-2, to=4-3]
        \arrow[from=4-2, to=5-2]
        \arrow[from=4-3, to=4-4]
        \arrow[from=4-3, to=5-3]
        \arrow[from=5-2, to=6-2]
        \arrow[from=5-3, to=6-3]
    \end{tikzcd}\]By \textit{(2)}, the second (and thus also the first) horizontal sequence is exact on the left. By snake lemma, $ \mathrm{Tor}^{A'}_1(M',A'/\mathfrak{p}') \cong \mathrm{Tor}^{A}_1(M,A/\mathfrak{p})=0 $ by \textit{(1)}.
\end{proof}

\begin{theorem}\label{B5}
    In the situation as beginning:\begin{enumerate}[(1)]
        \item There are three successive obstructions for the existence of an extension $ X' $ of $ X $ over $ C' $:\begin{itemize}
            \item $ \delta_1\in H^0(X_0,\mathcal{T}^2_{X_0}\otimes J) $ for the existence of local extension.
            \item $ \delta_2\in H^1(X_0,\mathcal{T}^1_{X_0}\otimes J) $ for the compatible isomorphism type of local extensions.
            \item $ \delta_3\in H^2(X_0,\mathcal{T}^0_{X_0}\otimes J) $ for the obstruction of gluing.
        \end{itemize}
        \item Fixed an extension $ X'_1 $ of $ X $ over $ C' $, there is an exact sequence\[
            0\longrightarrow H^1(X_0,\mathcal{T}^0_{X_0}\otimes J)\longrightarrow \mathrm{Def}(X/C,C')\longrightarrow H^0(X_0,\mathcal{T}^1_{X_0}\otimes J)\longrightarrow H^2(X_0,\mathcal{T}^0_{X_0}\otimes J).
        \]
        \item The automorphism group of $ X'_1 $ as an extension of $ X $ over $ C' $ is naturally isomorphic to the group $ H^0(X_0,\mathcal{T}_{X_0}\otimes J) $.
    \end{enumerate}
\end{theorem}

\begin{proof}
    For each open affine subscheme $ U_i\subseteq X $, by (\ref{B3}\textit{(1)}), there is an obstruction lying in $ H^0(U_i,\mathcal{T}^2_{U_i}\otimes J) $ for the existence of an extension $ U'_i $ of $ U_i $ over $ C' $. Since the definition of obstruction in (\ref{B3}\textit{(1)}) is compatible with localization, we may glue the data to a global obstruction $ \delta_1\in H^0(X_0,\mathcal{T}^2_{X_0}\otimes J) $.

    If $ \delta_1=0 $, then for each $ U_i $, there exists an extension $ U'_i $ of $ U_i $ over $ C' $. For each $ U_{ij}=U_i\cap U_j $, we have two deformations $ U'_i|_{U_{ij}} $ and $ U'_j|_{U_{ij}} $. By (\ref{B3}\textit{(2)}), their difference gives an element in $ H^0(U_{ij},\mathcal{T}^1\otimes J) $. The difference of three of these is zero on $ U_{ijk} $, so we get the second obstruction $ \delta_2\in H^1(X_0,\mathcal{T}_{X_0}^1\otimes J) $.

    If $ \delta_2=0 $, then there exists isomorphism $ \varphi_{ij}:U'_i|_{U_{ij}}\stackrel{\sim}{\longrightarrow}U'_j|_{U_{ij}} $ as the extension of $ U_{ij} $ over $ C' $ for each $ i,j $. Then $ \varphi_{ij}\circ\varphi_{jk}\circ\varphi_{ki}|_{U_{ijk}} $ gives an automorphism of $ U'_i|_{U_{ijk}} $. By (\ref{B4}), this gives an element in $ H^0(U_{ijk},\mathcal{T}^0\otimes J) $, and hence glue to the third obstruction $ \delta_3\in H^2(X_0,\mathcal{T}^0_{X_0}\otimes J) $.

    If $ \delta_3=0 $, then we can glue the extensions $ U'_i $ to get a global extension $ X' $ by gluing lemma.

    For \textit{(2)}, fixed an extension $ X'_1 $ of $ X $ over $ C' $. If $ X'_2 $ is another extension, by (\ref{B3}), their difference on each open affine subscheme $ U_i $ gives an element of $ H^0(U_i,\mathcal{T}^1_{X_0}\otimes J) $. These glue to a global element of $ H^0(X_0,\mathcal{T}^1_{X_0}\otimes J) $. Conversely, a global element of $ H^0(X_0,\mathcal{T}^1_{X_0}\otimes J) $ will give extensions of $ U_i $ that are isomorphic on the intersection $ U_{ij} $, and will there is an obstruction in $ H^2(X_0,\mathcal{T}^0_{X_0}\otimes J) $ to glue these to a global extension.

    Suppose that two extensions $ X'_2 $ and $ X'_3 $ give the same element in $ H^0(X_0,\mathcal{T}^1_{X_0}\otimes J) $. Then there is an isomorphism $ \phi_i:X'_2|_{U_i}\stackrel{\sim}{\longrightarrow}X'_3|_{U_i} $ as extension of $ U_i $. Then $ \psi_{ij}:=\varphi_j^{-1}\circ\varphi_i $ is an automorphism of $ X'_2|_{U_{ij}} $, which defines a element in $ H^0(U_{ij},\mathcal{T}^0_{X_0}\otimes J) $. These elements agree on $ U_{ijk} $, so we get an element of $ H^1(X_0,\mathcal{T}^0_{X_0}\otimes J) $. The vanishing of this element is equivalent to the existence of automorphism $ \theta_i:X'_2|_{U_i}\rightarrow X'_2|_{U_i} $ as the extension of $ U_i $ such that $ \psi_{ij}=\theta_j^{-1}\circ\theta_i $. Then $ \varphi_i\circ\theta_i $ glue to an isomorphism between $ X_2' $ and $ X_3' $ as the extension of $ X $. Hence we get the desired exact sequence for $ \mathrm{Def}(X/C,C') $.

    For \textit{(3)}, the automorphism group of $ X'_1|_{\mathrm{Spec}\, B} $ is naturally isomorphic to\[
        T^0(B_0/k,B_0\otimes J)=\mathrm{Hom}_k(\varOmega_{B_0/k},B_0\otimes J)=T_{B_0/k}\otimes J=(\mathcal{T}^0_{X_0}\otimes J)(\mathrm{Spec}\, B).
    \]Hence, the automorphism group of $ X'_1 $ is naturally isomorphic to $ H^0(X_0,\mathcal{T}^0_{X_0}\otimes J) $ by gluing.
\end{proof}

\begin{remark}\label{B6}
    The assumption $ \mathfrak{m}_{C'}J=0 $ is too strong for the discussion in the future. In fact, in most situations we only assume that $ J^2=0 $. The statements of the theorems above can be modified accordingly as follows:\begin{enumerate}[(1)]
        \item In the proof of (\ref{B2}), we show that the embedding extension of $ Y $ over $ C' $ in $ X' $ form a\[
            \mathrm{Hom}_B(I,J\otimes_C A)=\mathrm{Hom}_A(I/I^2,J\otimes_C A)=H^0(\mathrm{Spec}\, A,\mathcal{N}_{Y/X}\otimes f^*\tilde{J})
        \]pseudotorsor, where $ f:Y\rightarrow \mathrm{Spec}\, C $ is the structure map.
        \item In the proof of (\ref{B3}), the obstruction live in\[
            T^2(B/C,B_0\otimes_k J)=T^2(B/C,B\otimes_C J)=H^0(\mathrm{Spec}\, B,\mathcal{T}^2(\mathrm{Spec}\, B/\mathrm{Spec}\, C,f^*\tilde{J}))
        \]and we show that $ \mathrm{Def}(B/C,C') $ form a\[
            T^1(B/C,B\otimes_C J)=H^0(\mathrm{Spec}\, B,\mathcal{T}^1(\mathrm{Spec}\, B/\mathrm{Spec}\, C,f^*\tilde{J}))
        \]pseudotorsor, where $ f:\mathrm{Spec}\, B\rightarrow \mathrm{Spec}\, C $ is the structure map.
        \item Finally, (\ref{B5}) may be reformulated by replacing each cohomology group of the form $ H^q(X_0,\mathcal{T}^p_{X_0}\otimes J) $ with\[
            H^q(X,\mathcal{T}^p(X/\mathrm{Spec}\, C,f^*\tilde{J})),
        \]where $ f:X\rightarrow \mathrm{Spec}\, C $ is the structure map. Or more concisely, we may deonte it by $ H^q(X,\mathcal{T}^p(X/C,f^*\tilde{J})) $.
    \end{enumerate}
\end{remark}

\begin{corollary}\label{B7}
    If $ X_0 $ is nonsingular, then\begin{itemize}
        \item There is just one obstruction in $ H^2(X_0,\mathcal{T}_{X_0}\otimes J) $ for the existence of an extension $ X' $ of $ X $ over $ C' $.
        \item If such extension exist, their equivalence classes form a $ H^1(X_0,\mathcal{T}_{X_0}\otimes J) $-torsor.
    \end{itemize}
\end{corollary}

\begin{proof}
    By (\ref{A17}), the sheaves $ \mathcal{T}^1_{X_0} $ and $ \mathcal{T}^2_{X_0} $ are zero. Then the assertion follows from (\ref{B5}).
\end{proof}

\subsection{Functors of Artin rings}

Now we begin to formalize the idea of studying local deformations of a fixed object.

Still let $ k $ be a fixed algebraically closed ground field, and let $ \mathsf{C} $ be the category of local artinian $ k $-algebras with residue field $ k $. Let $ R $ be a complete $ k $-algebra, and for each $ A\in \mathsf{C} $, let $ h_{R}(A) $ be the set of $ k $-algebra homomorphisms $ \mathrm{Hom}(R,A) $. For any morphism $ A\rightarrow B $ in $ \mathsf{C} $ we get a map of sets $ h_R (A)\rightarrow h_R (B) $, so $ h_R $ is a covariant functor from $ \mathsf{C} $ to $ \mathsf{Set} $.

\begin{definition}
    A covariant functor $ F:\mathsf{C}\rightarrow \mathsf{Set} $ that is (naturally) isomorphic to a functor of the form $ h_R $ for some complete local $ k $-algebra $ R $ is called \textbf{pro-presentable}.
\end{definition}

To explain the nature of an isomorphism between $ h_R $ and $ F $, let us consider more generally any homomorphism of functors $ \varphi:h_R \rightarrow F $ for a complete local $ k $-algebra $ R $ with maximal ideal $ \mathfrak{m} $. In particular, for each $ n $ this will give a map $ \varphi_n:\mathrm{Hom}(R,R/\mathfrak{m}^n)\rightarrow F(R/\mathfrak{m}^n) $, and the image of the quotient map of $ R $ to $ R/\mathfrak{m}^n $ gives an element $ \xi_n\in F(R/\mathfrak{m}^n) $. These elements $ \xi_n $ are compatible, in the sense that the natural map $ R/\mathfrak{m}^{n+1}\rightarrow R/\mathfrak{m}^{n} $ induces a map of sets $ F(R/\mathfrak{m}^{n+1})\rightarrow F(R/\mathfrak{m}^{n}) $ that sends $ \xi_{n+1} $ to $ \xi_n $. Thus the collection $ \{\xi_n\} $ defines an element $ \xi\in \underleftarrow{\lim}\ F(R/\mathfrak{m}^n) $. We will call such a collection $ \xi=\{\xi_n\} $ a \textbf{formal family} of $ F $ over the ring $ R $.

Here it is useful to introduce the category $ \hat{\mathsf{C}} $ of complete local $ k $-algebras with residue field $ k $. The category $ \hat{\mathsf{C}} $ contains the category $ \mathsf{C} $, and we can extend any functor F on $ \mathsf{C} $ to a fucntor $ \hat{F} $ from $ \hat{\mathsf{C}} $ to $ \mathsf{Set} $ by defining $ \hat{F}(R)=\underleftarrow{\lim}\ F(R/\mathfrak{m}^n) $ for any $ R\in \hat{\mathsf{C}} $. In this notation, $ \hat{F}(R) $ is the set of formal families of $ F $ over $ R $.

Coversely, a formal family $ \xi=\{\xi_n\} $ of $ \hat{F}(R) $ defines a homomorphism of functors $ \varphi:h_R\rightarrow F $ as follows. For any $ A\in \mathsf{C} $ and any homomorphism $ f:R\rightarrow A $, since $ A $ is artinian, it factors through $ R/\mathfrak{m}^n $ for some $ n $, say $ f=g\pi $, where $ \pi: R\rightarrow R/\mathfrak{m}^n $ and $ g:R/\mathfrak{m}^n\rightarrow A $. Then let $ \varphi(f) $ be the image of $ \xi_n $ under the map $ F(g):F(R/\mathfrak{m}^n)\rightarrow F(A) $. It is easy to check that these constructions are well-defined and inverse to each other, so we have the following:

\begin{proposition}
    If $ F $ is a functor from $ \mathsf{C} $ to $ \mathsf{Set} $ and $ R $ is a complete local $ k $-algebra with residue field $ k $, then there is a natural bijection between the set $ \hat{F}(R) $ of formal families $ \{\xi_n\mid \xi_n \in F(R/\mathfrak{m}^n)\} $ and the set of homomorphisms of functors $ h_R $ to $ F $.
\end{proposition}

Thus, if $ F $ is pro-representable, there is an isomorphism $ \xi:h_R \rightarrow F $ for some $ R $, and we can think of $ \xi $ as an element of $ \hat{F}(R) $. We say that the pair $ (R,\xi) $ \textbf{pro-repsents} the functor $ F $. One can verify easily that if F is pro-representable, the pair $ (R,\xi) $ is unique up to isomorphism.

In many cases of interest, the functors we consider will not be pro-representable, so we define the weaker notions of having a versal family or a miniversal family.

\begin{definition}
    Let $ F:\mathsf{C}\rightarrow \mathsf{Set} $ be a functor. A pair $ (R,\xi) $ with $ R\in \hat{\mathsf{C}} $ and $ \xi\in \hat{F}(R) $ is a \textbf{versal family} for $ F $ if the associated map $ h_R \rightarrow F $ is strongly surjective. Here we say that a morphism of functors $ G\rightarrow F $ is \textbf{strongly surjective} if for every $ A\in\mathsf{C} $, the map $ G(A)\rightarrow F(A) $ is surjective, and furthermore, for every surjection $ B\rightarrow A $ in $ \mathsf{C} $, the map $ G(B)\rightarrow G(A)\times_{F(A)}F(B) $ is also surjective. In our case this means that given a map $ R\rightarrow A $ inducing an element $ \eta\in F(A) $, and given $ \theta\in F(B) $ mapping to $ \eta $, one can lift the map $ R\rightarrow A $ to a map $ R\rightarrow B $ inducing $ \theta $.

    If in addition the map $ h_R(D) \rightarrow F(D) $ is bijective, where $ D=k[t]/t^2 $ is the ring of dual numbers, we say that $ (R,\xi) $ is a \textbf{miniversal family}, or that the functor has a \textbf{pro-representable hull} $ (R,\xi) $. We say that $ (R,\xi) $ is a \textbf{universal family} if it pro-represents the functor $ F $.
\end{definition}

The following proposition explains the significance of this terminology.

\begin{proposition}
    Let $ (R,\xi) $ be a formal family of the functor $ F $.
    \begin{enumerate}[(1)]
        \item If $ (R,\xi) $ is a versal family, then for any other formal family $ (S,\eta) $, there is a ring homomorphism $ f:R\rightarrow S $ such that the induced map $ \hat{F}(R)\rightarrow\hat{F}(S) $ sends $ \xi $ to $ \eta $.
        \item If $ (R,\xi) $ is miniversal, then for any $ (S,\eta) $ the map $ f:R\rightarrow S $ of (a) induces a unique homomorphism $ R/\mathfrak{m}_R^2\rightarrow S/\mathfrak{m}_S^2 $.
        \item If $ (R,\xi) $ is a universal family, then for any $ (S,\eta) $ as in (a), the corresponding map $ f:R\rightarrow S $ is unique.
    \end{enumerate}
\end{proposition}

\begin{proof}
    \textit{(1)} Let $ (R,\xi) $ be a versal family, and $ (S,\eta) $ be any formal family. Then by definition we have a strongly surjective morphism of functor $ \varphi:h_R \rightarrow F $, determined by $ \xi $. For each $ n $, we have an element $ \eta_n\in F(S/\mathfrak{m}^n_S) $. Since $ \varphi $ is surjective, we can lift it to an element $ \theta_n\in h_R (S/\mathfrak{m}_S^n) $, i.e., a homomorphism $ \theta_n:R\rightarrow S/\mathfrak{m}^n_S $ such that the induced map $ F(R/\mathfrak{m}^n_R)\rightarrow F(S/\mathfrak{m}^n_S) $ sends $ \xi_n $ to $ \eta_n $. Furthermore, because of the strong surjectivity of $ \varphi $, starting from $ \theta_1 $, we can lift successively and obtain a compatible family $ \{\theta_n\} $, that is, so that each $ \theta_n $ is obtained by $ \theta_{n+1} $ followed by the natural map $ S/\mathfrak{m}^{n+1}_S\rightarrow S/\mathfrak{m}^n_S $. Thus the maps $ \theta_n $ determine a homomorphism $ R\rightarrow \underleftarrow{\lim}S/\mathfrak{m}^n_S $, which is equal to $ S $, since $ S $ is a complete local ring. This map $ f:R\rightarrow S $ then induces the map $ \hat{F}(R)\rightarrow \hat{F}(S) $ that sends $ \xi $ to $ \eta $ by construction.

    \textit{(2)} If $ (R,\xi) $ is miniversal, to show that the induced map $ \bar{f}:R/\mathfrak{m}^2_R\rightarrow S/\mathfrak{m}^2_S $ is unique, it is sufficient to show that for every map $ S/\mathfrak{m}^2_S\rightarrow D=k[t]/t^2 $, the induced map $ R/\mathfrak{m}^2_R\rightarrow D $ is unique. But since $ h_R(D)\rightarrow F(D) $ is bijective, by definition, this follows immediately.

    \textit{(3)} If $ (R,\xi) $ is a universal family, then $ h_R\rightarrow F $ is an isomorphism, so $ \eta $ determines a unique map $ h_S\rightarrow h_R $, and hence a unique $ f:R\rightarrow S $.
\end{proof}

\subsection{Schlessinger's criterion}

In this section we will prove Schlessinger's theorem, which gives criteria for a functor of Artin rings to have a versal family, a miniversal family, or to be pro-representable.

We keep the notation of the previous subsection: $ k $ is a fixed algebraically closed field, $ \mathsf{C} $ is the category of local artinian $ k $-algebras, and $ F $ is a covariant functor form $ \mathsf{C} $ to $ \mathsf{Set} $. Note that the category $ \mathsf{C} $ has fibered direct products. If $ A'\rightarrow A $ and $ A''\rightarrow A $ are morphisms in $ \mathsf{C} $, we take $ A'\times_A A'' $ to be the set-theoretic fibered product $ \{(a',a'')\mid a'\ \text{and}\ a''\ \text{have the same image in}\ A\} $. The ring operations extend naturally, giving another object of $ \mathsf{C} $.

It is convenient to introduce the notation $ t_F $ for $ F(D) $, where $ D $ is the ring of dual numbers. We call this the \textbf{tangent space} of $ F $. Similarly $ t_R $ denotes the tangent space of the functor $ h_R $, which is just $ \mathrm{Hom}_k (R,D) $, and is equal to the dual vector space of $ \mathfrak{m}_R/\mathfrak{m}_R^2 $.

A \textbf{small extension} in $ \mathsf{C} $ is a surjective map $ A'\rightarrow A $ whose kernel $ I $ is a one-dimensional $ k $-vector space.

\begin{proposition}\label{schne}
    If $ F $ has a versal family, then:
    \begin{enumerate}[(1)]
        \item $ F(k) $ has just one element.
        \item For any morphisms $ A'\rightarrow A $ and $ A''\rightarrow A $ in $ \mathsf{C} $, the natural map\[
            F(A'\times_A A'')\rightarrow F(A')\times_{F(A)}F(A'')
        \]is surjective.
        
        If furthermore $ F $ has a miniversal family, then:
        \item For any $ A\in \mathsf{C} $, considering the maps $ A\rightarrow k $ and $ D\rightarrow k $, the map of \textit{(b)} above,\[
            F(A\times_k D)\rightarrow F(A)\times_{F(k)}F(D)
        \]is bijective.\label{hhh}
        \item $ F(D)=t_F $ has a natural structure of a finite-dimensional $ k $-vector space.
        \item For any small extension $ p:A'\rightarrow A $ and any element $ \eta\in F(A) $, there is a transitive group action of the vector space $ t_F $ on the set $ p^{-1}(\eta) $ (provided it is nonempty).\label{ovo}
        
        Finally, if $ F $ is pro-representable, then:
        \item The maps of \textit{(b)} are all bijective and the action of \textit{(e)} is bijective whenever $ p^{-1}(\eta) $ is nonempty.
    \end{enumerate}
\end{proposition}

\begin{proof}
    \textit{(1)} Since $ \mathrm{Hom}(R,k)\rightarrow F(k) $ is surjective, and $ \mathrm{Hom}(R,k) $ has just one element, so does $ F(k) $.

    \textit{(2)} Given elements $ \eta'\in F(A') $ and $ \eta''\in F(A'') $ mapping to the same element $ \eta\in F(A) $, by the strong surjective property of a versal family, there are compatible homomorphisms of $ R $ to $ A' $, $ A $, and $ A'' $ inducing these elements. Then there is a unique map of $ R $ to the product $ A'\times_{A}A'' $ inducing the given maps of $ R $ to $ A' $ and $ A'' $. This in turn defines an element of $ F(A'\times_{A}A'') $ that restricts to $ \eta' $ and $ \eta'' $ as required.

    \textit{(3)} Suppose we are given $ \eta\in F(A) $ and $ \xi\in F(D) $. We know from \textit{(2)} that there are elements of $ F(A\times_k D) $ lying over the pair $ (\eta,\xi) $. Suppose $ \theta_1 $ and $ \theta_2 $ are two such. Choose a homomorphism $ u:R\rightarrow A $ inducing $ \eta $. By the strong surjectivity we can lift $ u $ to $ v_1 $ and $ v_2:R\rightarrow A[t]/t^2 $, inducing $ \theta_1 $ and $ \theta_2 $. Since $ \theta_1 $ and $ \theta_2 $ both lie over $ \xi $, the projections of $ v_1 $ and $ v_2 $ to $ D $ both induce $ \xi $. By the hypothesis of miniversality, $ t_R\rightarrow t_F $ is bijective, so these restrictions are equal. Since $ v_1 $ and $ v_2 $ also induce the same map $ u:R\rightarrow A $, we obtain $ v_1=v_2 $, and hence $ \theta_1=\theta_2 $.

    \textit{(4)} By miniversality, $ t_R\rightarrow t_F $ is bijective, so we can just carry over the vector space structure on $ t_R $ to $ t_F $. But this structure can also be recovered intrinsically, using only the functorial properties of $ F $ and \textit{(3)} above (\ref{vcst}).

    \textit{(5)} Let $ A'\rightarrow A $ be a small extension with kernel $ I\cong k $. Note that $ A'\times_A A'\cong A'\times_k k[I] $ by sending $ (x,y)\mapsto (x,x_0 +y-x) $,where $ x_0\in k $ is the residue of $ x $ mod $ \mathfrak{m} $. Consider the surjective map\[
        F(A'\times_A A')\rightarrow F(A')\times_{F(A)}F(A')
    \]of \textit{(2)}. Using the isomorphism above and \textit{(3)}, we can reinterpret the left-hand side as $ F(A'\times_k k[I])\cong F(A')\times t_F $, and we get a surjective map\[
        F(A')\times t_F \rightarrow F(A')\times_{F(A)}F(A')
    \]that is an isomorphism on the first factor. If we take $ \eta\in F(A) $ and fix $ \eta'\in p^{-1}(\eta) $ then we get a surjective map\[
        \{\eta'\}\times t_F\rightarrow \{\eta'\}\times p^{-1}(\eta),
    \]and this gives a transitive group action of $ t_F $ on $ p^{-1}(\eta) $.

    \textit{(6)} If $ F $ is pro-representable the proof of \textit{(2)} shows that these maps are all bijective. It follows that the action of \textit{(5)} is also bijective.
\end{proof}

\begin{example}\label{vcst}
    Let $ F:\mathsf{C}\rightarrow\mathsf{Set} $ be a functor having properties \textit{(1)} and \textit{(3)} of (\ref{schne}). Then $ t_F=F(D) $ has a natural structure of a $ k $-vector space as follows. For each $ \lambda\in k $, the ring homomorphism $ \lambda:D\rightarrow D,\ t\mapsto\lambda t $ induces a map $ \lambda_*:t_F\rightarrow t_F $. Taking two copies of $ D $, say $ D_1=k[t_1]/t_1^2 $ and $ D_2=k[t_2]/t_2^2 $, consider the ring homomorphism $ D_1 \times_k D_2\rightarrow D,\ t_i\mapsto t,\ i=1,2 $. This induces a map $ F(D_1\times_k D_2)\rightarrow F(D) $. Using \textit{(1)} and \textit{(3)} this gives a map $ \mu:t_F\times t_F\rightarrow t_F $. Then taking $ \lambda_* $ and $ \mu $ as scalar multiplication and addition makes $ t_F $ into a $ k $-vector space.
\end{example}

\begin{lemma}[Chevalley's theorem]\label{cheth}
    Let $ (A,\mathfrak{m}) $ be a complete local ring. Let $ \{\mathfrak{a}_n\} $ be a descending sequence of ideals, with $ \cap\mathfrak{a}_n=(0) $. Then for each $ s $, there is an $ n(s) $ such that for $ n\geq n(s) $, $ \mathfrak{a}_n\subseteq \mathfrak{m}^s $.
\end{lemma}

\begin{proof}
    \cite[Vol. \uppercase\expandafter{\romannumeral2}, \uppercase\expandafter{\romannumeral5}\uppercase\expandafter{\romannumeral3}, Thm. 13, p.270]{Zariski}.
\end{proof}

\begin{theorem}[Schlessinger's criterion]\label{jkl}
    The functor $ F:\mathsf{C}\rightarrow\mathsf{Set} $ has a miniversal family if and only if:
    \begin{enumerate}[]
        \item $ (H_0) $ $ F(k) $ has just one element.
        \item $ (H_1) $ $ F(A'\times_A A'')\rightarrow F(A')\times_{F(A)}F(A'') $ is surjective for every small extension $ A''\rightarrow A $.
        \item $ (H_2) $ The map of $ H_1 $ is bijective when $ A''=D $ and $ A=k $.
        \item $ (H_3) $ $ t_F $ is a finite-dimensional $ k $-vector space.
        
        Furthermore, $ F $ is pro-representable if and only if in addition:
        \item $ (H_4) $ For every small extension $ p:A''\rightarrow A $ and every $ \eta\in F(A) $ for which $ p^{-1}(\eta) $ is nonempty, the group action of $ t_F $ on $ p^{-1}(\eta) $ is bijective.
    \end{enumerate}
\end{theorem}

\begin{proof}
    The necessity has been proved.
    So now let $ F $ be a functor satisfying conditions $ H_0,H_1,H_2,H_3 $. First we will construct a ring $ R $ and a morphism $ h_R\rightarrow F $. Then we will show that it has the miniversal family property.

    We will define $ R $ and the map $ h_R\rightarrow F $ as an inverse limit of rings $ R_q $ and elements $ \xi_q\in F(R_q) $ that we construct inductively for $ q \geq 0 $. We take $ R_0=k $. Note that the vector space structure on $ t_F $ is already determined by conditions $ H_0,H_1,H_2 $ (\ref{vcst}), so that it makes sense to say in $ H_3 $ that it is finite-dimensional. Let $ x_1,\dots,x_r $ be a basis of the dual vector space $ t_F^* $, let $ S $ be the formal power series ring $ k[[x_1,\dots,x_r]] $, with maximal ideal $ \mathfrak{m} $, and take $ R_1=S/\mathfrak{m}^2 $. Then $ t_{R_1}\cong t_F $ by construction. Furthermore, by iterating the condition $ H_2 $, we obtain $ F(R_1)=F(k[x_1]\times\cdots\times k[x_r])\cong t_F\otimes_k t_F^* $. The natural element here gives $ \xi_1\in F(R_1) $, inducing the isomorphism $ t_{R_1}\cong t_F $.

    Now suppose we have constructed a compatible sequence $ (R_i,\xi_i) $ for $ i=1,\dots,q $, with $ \xi_i\in F(R_i) $, where $ R_i=S/J_i $, and $ \mathfrak{m}^{i+1}\subseteq J_i \subseteq J_{i-1} $, and for each $ i $ the natural map $ R_i\rightarrow R_{i-1} $ sends $ \xi_i $ to $ \xi_{i-1} $. Then, to construct $ R_{q+1} $, we look at ideals $ J $ in $ S $, with $ \mathfrak{m}J_q \subseteq J\subseteq J_q $, and take $ J_{q+1} $ to be the minimal such ideal $ J $ with the property that $ \xi_q\in F(R_q) $ lifts to an element $ \xi'\in F(S/J) $. To show that there is a minimal such $ J $, it will be sufficient to show that if $ J $ and $ K $ are two such, then their intersection $ J\cap K $ is another one. By enlarging $ J $ or $ K $ we may assume without loss of generality that $ J+K=J_q $. In that case $ S/(J\cap K)=(S/J)\times_{(S/J_q)}(S/K) $. Now the existence of liftings of $ \xi_q $ over $ S/J $ and $ S/K $ implies by condition $ H_1 $ the existence of a lifting over $ S/J\cap K $. Then we take $ R_{q+1}=S/J_{q+1} $, and $ \xi_{q+1} $ any lifting of $ \xi_q $, which exists by construction.

    Thus we obtain a surjective system of rings $ R_q $ and compatible elements $ \xi_q\in F(R_q) $. Let $ J=\cap J_q $ and take $ R $ to be $ S/J $. Then $ R $ is a complete local ring and $ R_q=R/\bar{J_q} $, where $ \bar{J_q}=J_q/J $. Since $ J_q\supseteq \mathfrak{m}^{q+1} $ for each $ q $, and conversely, for each $ s $, some $ \bar{J}_q\subseteq \mathfrak{m}_R^s $ (\ref{cheth}), it follows that the ideals $ \bar{J}_q $ form a base for the $ \mathfrak{m}_R $-adic topology of $ R $. Therefore $ R=\underleftarrow{\lim}\, R_q $, and we define $ \xi=\underleftarrow{\lim}\, \xi_q\in \hat{F}(R) $. This gives the desired map $ h_R\rightarrow F $.

    I claim that $ (R,\xi) $ is a miniversal family for $ F $. Since $ t_R \cong t_F $ by construction, we have only to show for any surjective map $ A'\rightarrow A $ and any $ \eta'\in F(A') $ restricting to $ \eta\in F(A) $, and any map $ R\rightarrow A $ inducing $ \eta $, that there exists a lifting to a map $ R\rightarrow A' $ inducing $ \eta' $. Since any surjective map factors into a sequence of small extensions, it suffices to treat the case of a small extension $ A'\rightarrow A $.

    Let $ u:R\rightarrow A $ induce $ \eta $. It will be sufficient to show that $ u $ lifts to some map $ v: R\rightarrow A' $. For then $ v $ will induce an element $ \eta''\in F(A') $ lying over $ \eta $. Because of~\ref{ovo}, whose proof used only~\ref{hhh}, which is our $ H_2 $, there is an element of $ t_F $ sending $ \eta'' $ to $ \eta' $ by the group action. This same $ t_F=t_R $ acts on the set of $ v:R\rightarrow A' $ restricting to $ u $, so then we can adjust $ v $ to a homomorphism $ v':R\rightarrow A' $ inducing $ \eta' $.

    Thus it remains to show that for a small extension $ A'\rightarrow A $, the given map $ u:R\rightarrow A $ lifts to a map $ v:R\rightarrow A' $. Since $ A $ is an Artin ring, $ u $ factors through $ R_q $ for some $ q $. On the other hand, $ R $ is a quotient of the power series ring $ S $, and the map $ u $ will lift to a map of $ S $ into $ A' $. Thus we get a commutative diagram
    % https://quiver.taskfirst.cn/#q=WzAsNixbMCwwLCJTIl0sWzEsMCwiUl9xXFx0aW1lc19BIEEnIl0sWzIsMCwiQSciXSxbMCwxLCJSIl0sWzEsMSwiUl9xIl0sWzIsMSwiQSJdLFswLDEsInciXSxbMSwyXSxbMiw1LCJwIl0sWzQsNV0sWzAsM10sWzMsNF0sWzEsNCwicCciXV0=
    \[\begin{tikzcd}
        S & {R_q\times_A A'} & {A'} \\
        R & {R_q} & A
        \arrow["w", from=1-1, to=1-2]
        \arrow[from=1-1, to=2-1]
        \arrow[from=1-2, to=1-3]
        \arrow["{p'}", from=1-2, to=2-2]
        \arrow["p", from=1-3, to=2-3]
        \arrow[from=2-1, to=2-2]
        \arrow[from=2-2, to=2-3]
    \end{tikzcd}\]
    Note that $ p':R_q\times_A A'\rightarrow R_q $ is also a small extension. If this map has a section $ s:R_q\rightarrow R_q\times_A A' $, then using $ s $ and the second projection we get a map $ R_q\rightarrow A' $ lifting $ u $, and we are done.

    If $ p' $ does not have a section, then I claim that the map $ w $ is surjective. Indeed, if $ w $ is not surjective, then $ \mathrm{Im}\, w $ is a subring mapping surjectively to $ R_q $. The kernel of $ \mathrm{Im}\, w\rightarrow R_q $ must be strictly contained in $ I=\ker p' $, which is a one-dimensional vector space, so this kernel is zero, the map $ \mathrm{Im}\, w\rightarrow R_q $ is an isomorphism, and this gives a section. Contradiction!

    Thus knowing that $ w $ is surjective, let $ J=\ker w $. Then $ J\subseteq J_q $, since $ S $ maps to $ R_q $ via $ w $. On the other hand, $ J\supseteq \mathfrak{m}J_q $, since $ p' $ is a small extension. But also $ \xi_q\in F(R_q) $ and there is an $ \eta'\in F(A') $ lying over $ \eta\in F(A) $, so by $ H_1 $, there is a $ \xi'\in F(R_q\times_A A') $ lying over both of these. Since $ R_q\times_A A'=S/J $, this ideal $ J $ satisfies the condition imposed in the construction. Therefore $ J\supseteq J_{q+1} $ and $ w $ factors through $ R_{q+1} $. This gives the required lifting of $ R $ to $ A' $, and completes the proof that $ (R,\xi) $ is a miniversal family for $ F $.

    Finally, suppose in addition that $ F $ satisfies $ H_4 $. To show that $ h_R(A)\rightarrow F(A) $ is bijective for all $ A $, it will be sufficient to show inductively, strating $ A=k $, that for any small extension $ p:A'\rightarrow A $ we have that $ \mathrm{Hom}(R,A')\rightarrow \mathrm{Hom}(R,A)\times_{F(A)}F(A') $ is bijective. So fix $ u\in \mathrm{Hom}(R,A) $ and the corresponding $ \eta\in F(A) $. If there is no map of $ R $ to $ A' $ lying over $ u $, then there is also no $ \eta'\in F(A') $ lying over $ \eta $, and there is nothing to prove. On the other hand, if $ p^{-1}(\eta) $ is nonempty, then the action of $ t_F $ on $ p^{-1}(\eta) $ is bijective by $ H_4 $, and the action of $ t_R $ on the set of homomorphism $ R\rightarrow A' $ lying over $ u $ is bijective, since $ h_R $ is pro-representable, and $ t_R\cong t_F $ by miniversality, so our map is bijective as required.
\end{proof}

\subsection{Deformation functor}

We now discuss the question of pro-representability or existence of a miniversal family of the deformation functor of a scheme.

Let $ X_0 $ be a scheme over $ k $ and $ A $ be a Artin ring with residue field $ k $. Define the following \textbf{deformation fucntor}:
% \[
%     F_1 (A)=\{\text{isomorphism classes of flat families}\ X/A\ \text{such that}\ X\otimes_A k\ \text{is isomorphic to}\ X_0\},
% \]and
\[
    F(A):=\mathrm{Def}(X_0/k,A)
\]

% Comparing with $ F_1 $, $ F $ will leave possible automorphisms of $ X_0 $ out of the picture and thus simplify the discussion. We first investigate $ F $ and consider the relation between these two functors later.

\begin{theorem}\label{yui}
    Let $ X_0 $ be a scheme over $ k $. Then the functor $ F $ has a miniversal family under either of the two folloing hypotheses:
    \begin{enumerate}[(1)]
        \item $ X_0 $ is affine with isolated singularities and of finite type.
        \item $ X_0 $ is projective.
    \end{enumerate}
\end{theorem}

\begin{proof}
    We verify the conditions of Schlessinger's criterion (\ref{jkl}).

    ($ H_0 $) $ F(k) $ consists of the single object $ (X_0,\mathrm{id}) $. If $ \sigma $ is an automorphism of $ X_0 $, the object $ (X_0,\sigma) $ is isomorphic to $ (X_0,\mathrm{id}) $ by the map $ \sigma: X_0\rightarrow X_0 $.

    ($ H_1 $) Suppose we are given a small extension $ A''\rightarrow A $ and any map $ A'\rightarrow A $, and suppose we are given objects $ X'\in F(A') $, $ X''\in F(A'') $ restricting to $ X\in F(A) $. Then $ X'\otimes_{A'}A\cong X $, the isomorphism being compatible wtih the maps from $ X_0 $, so we can choose a closed immersion $ X\hookrightarrow X' $ inducing this isomorphism. Similarly choose $ X\hookrightarrow X'' $. Then we define $ X^* $ by the fibered product of sheaves of rings $ \mathcal{O}_{X^*}\mathcal{O}_{X'}\times_{\mathcal{O}_X}\mathcal{O}_{X''} $, and $ X^* $ will be an object of $ F(A^*) $ reducing to $ X' $ and $ X'' $, where $ A^* =A'\times_A A'' $.

    $ (H_2) $ Suppose $ A=k $ in the situation of $ (H_1) $ (which effectively means $ A''\cong k[t] $) and let $ X^* $ be constructed as in $ (H_1) $. If $ W $ is any other object of $ F(A^*) $ restricting to $ X' $ and $ X'' $ respectively, then we can choose immersions $ X'\hookleftarrow W $ and $ X''\hookrightarrow W $ inducing these isomorphisms. Since these maps are all compatible with the immersion from $ X_0 $, they agree with the chosen maps $ X\hookrightarrow X' $ and $ X\hookrightarrow X'' $, since in this case $ X=X_0 $. Now by the universal property of fibered product of rings, there is a map $ X^*\rightarrow W $ compatible with the above maps. Since $ X^* $ and $ W $ are both flat over $ A^* $, and the map becomes an isomorphism when restricted to $ X_0 $, we find that $ X^* $ is isomorphic to $ W $, and hence they are equal as elements of $ F(A^*) $.

    $ (H_3) $ Here is the only place in the proof that we need the hypothesis \textit{(1)} or \textit{(2)}.
    
    \quad \textit{(1)} Let $ X_0=\mathrm{Spec}\, B $. Then $ t_F = T_{B/k}^1 $ by (\ref{B3}). This module is supported at the finite number of singular points of $ X_0 $ (\ref{A17}), so has finite length (\ref{A10}), i.e., $ t_F $ is a finite-dimensional vector space.

    \quad \textit{(2)} For arbitrary $ X_0 $, the tangent space $ t_F $ corresponds to deformations over the ring of dual numbers $ D $. Because of the exact sequence (\ref{B5}\textit{(2)})\[
        0\longrightarrow H^1(X_0,\mathcal{T}_{X_0})\longrightarrow \mathrm{Def}(X_0/k,D)\longrightarrow H^0(X_0,\mathcal{T}_{X_0}^1)\longrightarrow\cdots,
    \]we see that if $ X_0 $ is projective, the two outside groups are finite-dimensional vector spaces, and so $ \mathrm{Def}(X_0/k,D) $ is also.

    Thus conditions $ (H_0)$--$ (H_3) $ are satisfied, and $ F $ has a miniversal family.
\end{proof}

Next, we consider conditions under which the functor $ F $ is actually pro-representable.

\begin{theorem}\label{ert}
    Let $ X_0/k $ be given and assume the hypotheses of (\ref{yui}) satisfied. Then the functor $ F $ of deformations of $ X_0 $ is pro-representable if and only if for each small extension $ A'\rightarrow A $, and for each deformation $ X' $ over $ A' $ restricting to a deformation $ X $ over $ A $, the natural map $ \mathrm{Aut}(X'/X_0)\rightarrow \mathrm{Aut}(X/X_0) $ of automorphisms of $ X' $ (and $ X $) restricting the identity automorphism of $ X_0 $ is surjective.
\end{theorem}

\begin{proof}
    Suppose that $ \mathrm{Aut}(X'/X_0)\rightarrow \mathrm{Aut}(X/X_0) $ is surjective for every $ X' $ lying over $ X $. We need to verify condition $ (H_4) $ of (\ref{jkl}). Let $ X'_1 ,X'_2 $ be elements of $ F(A') $ inducing the same element $ X $ of $ F(A) $. If $ X\hookrightarrow X'_1 $ and $ X\hookrightarrow X'_2 $ are maps inducing the isomorphisms $ X'_1\otimes_{A'}A\cong X $ and $ X'_2\otimes_{A'}A\cong X $, and if $ X'_1 $ and $ X'_2 $ are isomorphic as deformations of $ X_0 $, then I claim that the inclusions $ X\hookrightarrow X'_1 $ and $ X\hookrightarrow X'_2 $ are isomorphic as extensions of $ X $ over $ A' $.

    Indeed, let $ u':X'_1\rightarrow X'_2 $ be an isomorphism over $ X_0 $. Then $ u=u'\otimes_{A'}A $ is an automorphism of $ X $ over $ X_0 $. By hypothesis this lifts to an automorphism $ \sigma $ of $ X'_1 $. Then $ v=u\circ \sigma^{-1}:X'_1\rightarrow X'_2 $ is an isomorphism inducing the identity on $ X $, so $ X'_1 $ and $ X'_2 $ are equivalent as extensions of $ X $ over $ A' $.

    Now by (\ref{B5}\textit{(2)}), $ \mathrm{Def}(X/A,A') $ is a principal homogeneous space under the action of $ t_F $ ($ D\rightarrow k $ is also a small extension), so condition $ (H_4) $ of Schlessinger's criterion is satisfied, and $ F $ is pro-representable.

    Conversely, suppose that $ F $ is pro-representable, let $ X'\in F(A') $ restrict to $ X\in F(A) $, and choose a map $ u:X\hookrightarrow X' $ inducing the isomorphism $ X\stackrel{\sim}{\longrightarrow}X'\otimes_{A'}A $. Let $ \sigma\in \mathrm{Aut}(X/X_0) $. Then $ u'=u\circ \sigma:X\hookrightarrow X' $ gives another element of $ \mathrm{Def}(X/A,A') $, and so $ u $ and $ u' $ differ by an element of $ t_F $, by (\ref{B5}\textit{(2)}). But $ u $ and $ u' $ define the same element $ X'\in F(A') $, lying over $ X $, so by condition $ (H_4) $, this element of $ t_F $ must be zero. Hence $ u $ and $ u' $ are equal as elements of $ \mathrm{Def}(X/A,A') $; in other words there exists an isomorphism $ \tau:X'\rightarrow X' $ over $ X_0 $ such that $ u'=\tau\circ u $. Restricting to $ X $ we obtain $ \sigma=\tau|_X $. Thus $ \tau\in \mathrm{Aut}(X'/X_0) $ lifts $ \sigma $, and tha map is surjective.
\end{proof}

\begin{corollary}\label{rty}
    Let $ X_0/k $ be a projective scheme with $ H^0(X_0,\mathcal{T}_{X_0})=0 $. Then the functor of deformations of $ X_0/k $ is pro-representable.
\end{corollary}

\begin{proof}
    We will show, by induction on the length of $ A $, that for any deformation $ X $ of $ X_0 $ over $ A $, $ \mathrm{Aut}(X/X_0)=\{\mathrm{id}\} $. Then obviously the condition of (\ref{ert}) will be satisfied.

    We start the induction by noting that $ \mathrm{Aut}(X_0/X_0)=\{\mathrm{id}\} $. Inductively, assume that $ \mathrm{Aut}(X/X_0)=\{\mathrm{id}\} $, where $ X $ is a deformation over $ A $. Consider a small extension $ A'\rightarrow A $ and any $ X'\in F(A') $ restricting to $ X $. Choose a map $ X\hookrightarrow X' $ inducing the isomorphism $ X\stackrel{\sim}{\rightarrow}X'\otimes_{A'}A $. Any automorphism of $ X' $ restricts to the identity on $ X $, by the induction hypothesis, so it is an automorphism of the deformation $ X\hookrightarrow X' $. Since these are classfied by $ H^0(X_0,\mathcal{T}_{X_0})=0 $ (\ref{B5}\textit{(3)}), this automorphism is the identity.
\end{proof}

\begin{example}
    Let $ X_0 $ be a nonsingular projective curve over the algebraically closed field $ k $. If the genus is zero, $ X_0=\mathbb{P}^1_k $. In this case $ H^0(X_0,\mathcal{T}_{X_0}) $ has dimension 3, but still, it can be proved that the functor of deformation of $ X_0/k $ is pro-representable. Thus the condition of (\ref{rty}) is not necessary for the functor to be pro-representable.
\end{example}

\subsection{Ran-Kawamata $ T^1 $ lifting}

\begin{definition}\label{ccc}
    An \textbf{obstruction theory} for a functor $ F:\mathsf{C}\rightarrow\mathsf{Set} $ is a $ k $-vector space $ V $, together with, for every exact sequence $ 0\rightarrow J\rightarrow C'\rightarrow C\rightarrow 0 $ such that $ \mathfrak{m}_{C'}J=0 $, and for every $ u\in F(C) $, an element $ \varphi(u,C')\in V\otimes_k J $ such that\begin{itemize}
        \item $ \varphi(u,C')=0 $ if and only if $ u\in\mathrm{Im}\, F(C'\rightarrow C) $.
        \item $ \varphi $ is functorial in the sense that if% https://quiver.taskfirst.cn/#q=WzAsMTAsWzAsMCwiMCJdLFsxLDAsIkpfMSJdLFsyLDAsIkMnXzEiXSxbMywwLCJDXzEiXSxbNCwwLCIwIl0sWzAsMSwiMCJdLFsxLDEsIkpfMiJdLFsyLDEsIkNfMiciXSxbMywxLCJDXzIiXSxbNCwxLCIwIl0sWzAsMV0sWzEsMl0sWzIsM10sWzMsNF0sWzUsNl0sWzYsN10sWzcsOF0sWzgsOV0sWzEsNl0sWzIsN10sWzMsOF1d
    \[\begin{tikzcd}
        0 & {J_1} & {C'_1} & {C_1} & 0 \\
        0 & {J_2} & {C_2'} & {C_2} & 0
        \arrow[from=1-1, to=1-2]
        \arrow[from=1-2, to=1-3]
        \arrow[from=1-2, to=2-2]
        \arrow[from=1-3, to=1-4]
        \arrow[from=1-3, to=2-3]
        \arrow[from=1-4, to=1-5]
        \arrow[from=1-4, to=2-4]
        \arrow[from=2-1, to=2-2]
        \arrow[from=2-2, to=2-3]
        \arrow[from=2-3, to=2-4]
        \arrow[from=2-4, to=2-5]
    \end{tikzcd}\]is a commutative diagram such that $ \mathfrak{m}_{C'_i}J_i=0 $, then we have the following commutative diagram% https://quiver.taskfirst.cn/#q=WzAsNixbMCwwLCJGKEMnXzEpIl0sWzEsMCwiRihDXzEpIl0sWzIsMCwiVlxcb3RpbWVzX2tKXzEiXSxbMCwxLCJGKEMnXzIpIl0sWzEsMSwiRihDXzIpIl0sWzIsMSwiVlxcb3RpbWVzX2tKXzIiXSxbMCwzXSxbMCwxXSxbMSwyXSxbMiw1XSxbMSw0XSxbNCw1XSxbMyw0XV0=
    \[\begin{tikzcd}
        {F(C'_1)} & {F(C_1)} & {V\otimes_kJ_1} \\
        {F(C'_2)} & {F(C_2)} & {V\otimes_kJ_2}.
        \arrow[from=1-1, to=1-2]
        \arrow[from=1-1, to=2-1]
        \arrow[from=1-2, to=1-3]
        \arrow[from=1-2, to=2-2]
        \arrow[from=1-3, to=2-3]
        \arrow[from=2-1, to=2-2]
        \arrow[from=2-2, to=2-3]
    \end{tikzcd}\]
    \end{itemize}
\end{definition}

\begin{corollary}\label{B8}
    The deformation functor of smooth scheme $ X $ over $ k $ has obstruction theory.
\end{corollary}

\begin{proof}
    It follows from (\ref{B7}) and take $ V=H^2(X,\mathcal{T}_X) $.
\end{proof}

Now let $ k $ be a algebraically closed field with $ \mathrm{char}\, k=0 $ and $ F:\mathsf{C}\rightarrow\mathsf{Set} $ still be a deformation functor. We also assume that the obstruction can be calculated by a $ k $-vector space $ T^2 $ (\ref{ccc}).

For $ n\in \mathbb{Z}_{\geq 0} $, let\[
    A_n=k[t]/(t^{n+1}),\ B_n=A_n\otimes A_1=k[x,y]/(x^{n+1},y^2),\ C_n=k[x,y]/(x^{n+1},x^n y,y^2).
\]We have natural homomorphism $ \alpha_n:A_{n+1}\rightarrow A_n $, $ \beta_n:B_n\rightarrow A_n $, and $ \gamma_n:B_n\rightarrow C_n $.

For $ X_n\in F(A_n) $, let\[
    T^1(X_n/A_n)=\{Y_n\in F(B_n)\mid F(\beta_n)(Y_n)=X_n\}.
\]

\begin{theorem}[\cite{Kawamata92}]
    Assume that $ F $ is pro-representable by $ (R,\xi) $, and that the following $ T^1 $ lifting property holds: for any positive integer $ n $ and any $ X_{n+1}\in F(A_{n+1}) $, the natural homomorphism\[
        T^1(\alpha_n):T^1(X_{n+1}/A_{n+1})\rightarrow T^1(X_n/A_n)
    \]is surjective for $ X_n=F(\alpha_n)(X_{n+1}) $. Then $ F $ is unobstructed.
\end{theorem}

\begin{proof}
    It suffices to show that $ F(\alpha_n) $ is surjective for all $ n $. Indeed, according to the proof of Schlessinger's criterion, we have $ R=k[[x_1,\dots,x_r]]/J $, where $ J\subseteq\mathfrak{m}^2_R $ and $ r=\dim t_F $. Since $ J\subseteq\mathfrak{m}^2_R $, we can define the morphism $ g_2:R\rightarrow A_1 $ by $ x_i\mapsto a_i t $ for any $ a_i\in k $. Since $ F(\alpha_n) $ is surjective for all $ n $, there is a lifting $ g_n:R\rightarrow A_n $ of $ g_{n-1} $ by induction for all $ n\geq 3 $. This define a morphism $ g:R\rightarrow k[[t]] $ by $ x_i\mapsto a_i t $. Suppose that $ J $ contains a nonzero power series $ f=f_d+f_{d+1}+\cdots $, where $ f_n $ is the homogeneous polynomial of degree $ n $, $ f_d\neq 0 $ and $ d\neq 0 $. Take $ a_i $ such that $ f_d(a_1,\dots,a_r)\neq 0 $, then\[
        g(f)=t^d f_d(a_1,\dots,a_r)+t^{d+1}f_{d+1}(a_1,\dots,a_r)+\cdots\neq 0
    \]leading a contradiction. Hence, $ R=k[[x_1,\dots,x_r]] $ is a power-series ring over $ k $, $ F=h_R $ is clearly unobstructed for lifting of any map in $ \mathsf{C} $.

    Consider the following commutative diagram% https://quiver.taskfirst.cn/#q=WzAsNixbMCwwLCJEKEFfe24rMX0pIl0sWzAsMSwiRChCX24pIl0sWzEsMSwiRChDX24pIl0sWzEsMCwiRChBX24pIl0sWzIsMCwiVF4yXFxvdGltZXModF57bisxfSkiXSxbMiwxLCJUXjJcXG90aW1lcyh4Xm55KSJdLFswLDMsIkQoXFxhbHBoYV9uKSJdLFszLDQsIlxcZGVsdGFfMSJdLFsyLDUsIlxcZGVsdGFfMiJdLFswLDEsIkQoXFx2YXJlcHNpbG9uX24pIl0sWzEsMiwiRChcXGdhbW1hX24pIl0sWzMsMiwiRChcXHZhcmVwc2lsb24nX24pIl0sWzQsNSwiXFx2YXJwaGkiXV0=
    \[\begin{tikzcd}
        {F(A_{n+1})} & {F(A_n)} & {T^2\otimes(t^{n+1})} \\
        {F(B_n)} & {F(C_n)} & {T^2\otimes(x^ny)}
        \arrow["{F(\alpha_n)}", from=1-1, to=1-2]
        \arrow["{F(\varepsilon_n)}", from=1-1, to=2-1]
        \arrow["{\delta_1}", from=1-2, to=1-3]
        \arrow["{F(\varepsilon'_n)}", from=1-2, to=2-2]
        \arrow["\varphi", from=1-3, to=2-3]
        \arrow["{F(\gamma_n)}", from=2-1, to=2-2]
        \arrow["{\delta_2}", from=2-2, to=2-3]
    \end{tikzcd}\]where $ \varepsilon_n(t)=x+y $ and $ \varepsilon'_n(t)=x+y $. By the definition of $ T^2 $, we have \[ \mathrm{Im}\, F(\alpha_n)=\{X_n\in F(A_n)\mid\delta_1(X_n)=0\},\ \mathrm{Im}\, F(\gamma_n)=\{Y_n\in F(C_n)\mid\delta_2(Y_n)=0\}. \]Notice that $ \varphi $ is an isomorphism, since $ \varphi:t^{n+1}\mapsto(n+1)x^n y $ and $ \mathrm{char}\, k=0 $. Hence, to show that $ F(\alpha_n) $ is surjective, it suffices to show that $ F(\gamma_n) $ is surjective.

    From the following diagram% https://quiver.taskfirst.cn/#q=WzAsNSxbMSwwLCJCX24iXSxbMCwxLCJCX3tuLTF9XFx0aW1lc197QV97bi0xfX1BX249Q19uIl0sWzIsMSwiQV9uIl0sWzAsMywiQl97bi0xfSJdLFsyLDMsIkFfe24tMX0iXSxbMCwxLCJcXGdhbW1hX24iLDJdLFswLDIsIlxcYmV0YV9uIl0sWzAsM10sWzEsM10sWzMsNCwiXFxiZXRhX3tuLTF9IiwyXSxbMiw0LCJcXGFscGhhX3tuLTF9Il1d
    \[\begin{tikzcd}
        & {B_n} & \\
        {B_{n-1}\times_{A_{n-1}}A_n=C_n} && {A_n} \\
        \\
        {B_{n-1}} && {A_{n-1}}
        \arrow["{\gamma_n}"', from=1-2, to=2-1]
        \arrow["{\beta_n}", from=1-2, to=2-3]
        \arrow[from=1-2, to=4-1]
        \arrow[from=2-1, to=4-1]
        \arrow["{\alpha_{n-1}}", from=2-3, to=4-3]
        \arrow["{\beta_{n-1}}"', from=4-1, to=4-3]
    \end{tikzcd}\]and the pro-representability of $ F $, we have\[
        F(B_n)\stackrel{F(\gamma_n)}{\longrightarrow}F(C_n)\cong F(B_{n-1})\times_{F(A_{n-1})}F(A_n).
    \]By the $ T^1 $ lifting property, for any $ Z_n\in F(C_n) $, there exists $ Y_n\in F(B_n) $ whose image in $ F(B_{n-1}) $ and $ F(A_n) $ coincide with image of $ Z_n $, respectively. Hence, $ F(\gamma_n)(Y_n)=Z_n $.
\end{proof}

\begin{remark}\label{qqq}
    In~\cite{Kawamata97}, Kawamata replaced the definition of $ T^1 $ space by $ T^1(X_n/A_n):=\mathrm{Def}(X_n/A_n,B_n) $. With this new definition, the above theorem holds without the assumption that $ F $ is pro-representable.

    As in the argument above, it suffices to show that the map $ F(\gamma_n):F(B_n)\rightarrow F(C_n) $ is surjective. For $ Z_n\in F(C_n) $, let\[
        Y_{n-1}=Z_n\times_{C_n}B_{n-1},\ X_n=Z_n\times_{C_n}A_n,\ X_{n-1}=Z_n\times_{C_n}A_{n-1}.
    \]By the $ T^1 $ lifting property, let $ [Y_n] $ be the equivalent classes in $ \mathrm{Def}(X_n/A_n,B_n) $ maps to $ [Y_{n-1}]\in\mathrm{Def}(X_{n-1}/A_{n-1},B_{n-1}) $. Then there exists an isomorphism $ \sigma:Y_{n-1}\rightarrow Y_{n}\times_{B_{n}}B_{n-1} $ compatible with the closed immersion of $ X_{n-1} $, and we get a diagram% https://quiver.taskfirst.cn/#q=WzAsNyxbMCwxLCJYX3tuLTF9Il0sWzEsMCwiWV97bi0xfSJdLFsxLDEsIlpfbiJdLFsxLDIsIlhfbiJdLFsyLDAsIllfblxcdGltZXNfe0Jfbn1CX3tuLTF9Il0sWzIsMSwiWV9uXFx0aW1lc197Ql9ufUNfbiJdLFszLDEsIllfbiJdLFswLDFdLFsxLDQsIlxcc2lnbWEiXSxbMSwyXSxbMCwzXSxbMCw0XSxbMywyXSxbMyw1XSxbMyw2XSxbNSw2XSxbNCw1XSxbNCw2XSxbNSwyLCI/IiwyLHsic3R5bGUiOnsiYm9keSI6eyJuYW1lIjoiZGFzaGVkIn19fV1d
    \[\begin{tikzcd}
        & {Y_{n-1}} & {Y_n\times_{B_n}B_{n-1}} & \\
        {X_{n-1}} & {Z_n} & {Y_n\times_{B_n}C_n} & {Y_n} \\
        & {X_n}
        \arrow["\sigma", from=1-2, to=1-3]
        \arrow[from=1-2, to=2-2]
        \arrow[from=1-3, to=2-3]
        \arrow[from=1-3, to=2-4]
        \arrow[from=2-1, to=1-2]
        \arrow[from=2-1, to=1-3]
        \arrow[from=2-1, to=3-2]
        \arrow[dashed, from=2-3, to=2-2]
        \arrow[from=2-3, to=2-4]
        \arrow[from=3-2, to=2-2]
        \arrow[from=3-2, to=2-3]
        \arrow[from=3-2, to=2-4]
    \end{tikzcd}\]By (\ref{tuo}), there exists an isomorphism $ Y_n\times_{B_n}C_n\stackrel{\sim}{\longrightarrow}Z_n $ compatible with the closed immersion of $ X_{n} $. This show that $ F(\gamma_n)(Y_n)=Z_n $.
\end{remark}

\begin{example}\label{iop}
    Suppose $ X_n $ is a smooth deformation of $ X $ over $ A_n $. By (\ref{B6}), the deformation space can be computed via the following exact sequence\[
        0\longrightarrow H^1(X_n,\mathcal{T}^0(X_n/A_n,f^*\tilde{J}))\longrightarrow\mathrm{Def}(X_n/A_n,B_n)\longrightarrow H^0(X_n,\mathcal{T}^1(X_n/A_n,f^*\tilde{J})),
    \]where $ f:X_n\rightarrow A_n $ is the structure morphism and $ J=\ker\beta_n $. Since $ \ker\beta_n $ is a free $ A_n $-module and $ X_n $ is smooth over $ A_n $, it follows from (\ref{A18}) that\[
        \mathcal{T}^i(X_n/A_n,f^*\tilde{J})=\mathcal{T}^i_{X_n/A_n}=\begin{cases}
            \mathcal{T}_{X_n/A_n}&\quad i=0,\\
            0&\quad i=1.
        \end{cases}
    \]So we conclude that $ \mathrm{Def}(X_n/A_n,B_n)=H^1(X_n,\mathcal{T}_{X_n/A_n}) $.
\end{example}

\begin{lemma}\label{ghu}
    Let $ A $ be a ring, $ J $ be a nilpotent ideal in $ A $, and $ u:M\rightarrow N $ be a homomorphism of $ A $-modules, with $ N $ falt over $ A $. If $ \bar{u}:M/JM\rightarrow N/JN $ is an isomorphism, then $ u $ is an isomorphism.
\end{lemma}

\begin{proof}
    Let $ K=\mathrm{coker}\, u $. Tensor the exact sequence\[
        M\rightarrow N\rightarrow K\rightarrow 0
    \]with $ A/J $, we have $ K/JK=0 $. Since $ J $ is nilpotent, we have $ K=0 $. If $ K'=\ker u $, then we get an exact sequence\[
        0\rightarrow K'/JK'\rightarrow M/JM\rightarrow N/JN\rightarrow 0
    \]by the flatness of $ N $. Hence, $ K'=0 $.
\end{proof}

\begin{lemma}\label{fff}
    Consider a commutative diagram% https://quiver.taskfirst.cn/#q=WzAsOCxbNCwxLCJNIl0sWzQsMywiQSJdLFszLDIsIkFfMiJdLFszLDAsIk1fMiJdLFswLDAsIk4iXSxbMCwyLCJCIl0sWzEsMywiQV8xIl0sWzEsMSwiTV8xIl0sWzMsMCwidV8yIl0sWzAsMSwiIiwwLHsic3R5bGUiOnsiaGVhZCI6eyJuYW1lIjoibm9uZSJ9fX1dLFszLDIsIiIsMCx7InN0eWxlIjp7ImhlYWQiOnsibmFtZSI6Im5vbmUifX19XSxbNCw3LCJwXzEiLDJdLFs3LDAsInVfMSJdLFs0LDMsInBfMiJdLFs0LDUsIiIsMCx7InN0eWxlIjp7ImhlYWQiOnsibmFtZSI6Im5vbmUifX19XSxbNyw2LCIiLDIseyJzdHlsZSI6eyJoZWFkIjp7Im5hbWUiOiJub25lIn19fV0sWzUsMl0sWzIsMV0sWzYsMV0sWzUsNl1d
    \[\begin{tikzcd}
        N &&& {M_2} & \\
        & {M_1} &&& M \\
        B &&& {A_2} \\
        & {A_1} &&& A
        \arrow["{p_2}", from=1-1, to=1-4]
        \arrow["{p_1}"', from=1-1, to=2-2]
        \arrow[no head, from=1-1, to=3-1]
        \arrow["{u_2}", from=1-4, to=2-5]
        \arrow[no head, from=1-4, to=3-4]
        \arrow["{u_1}", from=2-2, to=2-5]
        \arrow[no head, from=2-2, to=4-2]
        \arrow[no head, from=2-5, to=4-5]
        \arrow[from=3-1, to=3-4]
        \arrow[from=3-1, to=4-2]
        \arrow[from=3-4, to=4-5]
        \arrow[from=4-2, to=4-5]
    \end{tikzcd}\]of compatible ring and module homomorphisms, where $ B=A_1\times_A A_2 $, $ N=M_1\times_M M_2 $, and $ M_i $ is a flat $ A_i $-module for $ i=1,2 $. Suppose that $ A_1 $ is an artinian ring and\begin{enumerate}[(1)]
        \item $ A_2/J\stackrel{\sim}{\longrightarrow}A $, where $ J $ is an nilpotent ideal in $ A_2 $.
        \item $ u_i $ induces $ M_i\otimes_{A_i}A\stackrel{\sim}{\longrightarrow}M $.
    \end{enumerate}
    Then $ N $ is flat over $ B $, and $ p_i $ induces $ N\otimes_B A_i\stackrel{\sim}{\longrightarrow}M_i $.
\end{lemma}

\begin{proof}
    Let $ \{x_{1,i}\}_{i\in I}\subseteq M_1 $ such that $ \{\bar{x}_{1,i}\}_{i\in I} $ form a basis of $ M_1/\mathfrak{m}_{A_1}M_1 $. Then $ \bigoplus_{i\in I}A_1 \rightarrow M_1 $ defined by $ e_i\mapsto x_{1,i} $ is an isomorphism by (\ref{ghu}). Hence $ M_1 $ is a free $ A_1 $-module with basis $ \{x_{1,i}\}_{i\in I} $. By \textit{(2)}, $ M $ is the free module with basis $ \{u_1(x_{1,i})\}_{i\in I} $. Choosing $ x_{2,i}\in M_2 $ such that $ u_2(x_{2,i})=u_1(x_{1,i}) $. Then $ A_2 $-module homomorphism $ \sum A_2 x_{2,i}\rightarrow M_2 $ is isomorphism after modulo the ideal $ J $. Again, by (\ref{ghu}), $ M_2 $ is free module with basis $ \{x_{2,i}\}_{i\in I} $, and it is clear that $ N $ is free $ B $-module with basis $ \{x_{1,i}\times x_{2,i}\}_{i\in I} $. Then the projection $ p_i $ induces the isomorphism $ N\otimes_{B}A_i\stackrel{\sim}{\longrightarrow}M_i $.
\end{proof}

\begin{lemma}\label{tuo}
    With the notations in (\ref{fff}). Let $ L $ be a $ B $-module satisfies the commutative diagram% https://quiver.taskfirst.cn/#q=WzAsNCxbMCwwLCJMIl0sWzEsMCwiTV8yIl0sWzAsMSwiTV8xIl0sWzEsMSwiTSJdLFswLDEsInFfMiJdLFswLDIsInFfMSIsMl0sWzIsMywidV8xIiwyXSxbMSwzLCJ1XzIiXV0=
    \[\begin{tikzcd}
        L & {M_2} \\
        {M_1} & M
        \arrow["{q_2}", from=1-1, to=1-2]
        \arrow["{q_1}"', from=1-1, to=2-1]
        \arrow["{u_2}", from=1-2, to=2-2]
        \arrow["{u_1}", from=2-1, to=2-2]
    \end{tikzcd}\]where $ q_1 $ induces $ L\otimes_B A_1\stackrel{\sim}{\longrightarrow}M_1 $. Then the canonical morphism $ q_1\times q_2:L\rightarrow N=M_1\times_M M_2 $ is an isomorphism.
\end{lemma}

\begin{proof}
    Apply (\ref{ghu}) to the morphism $ u=q_1\times q_2 $. Note that $ B/(0\times J)\stackrel{\sim}{\longrightarrow}A_1 $.
\end{proof}

\subsection{Deformations of smooth Calabi-Yau varieties}

Before proving the unobstructedness of deformations of Calabi-Yau varieties, we need the following essential technique.

\begin{theorem}\label{kkk}
    Let $ f:X\rightarrow S=\mathrm{Spec}\, A_n $ be a projective and smooth morphism, where $ A_n=\mathbb{C} [t]/(t^{n+1}) $. Then $ H^q(X,\varOmega^p_{X/S}) $ is a free $ A_n $-module.
\end{theorem}

\begin{proof}
    Let $ A=A_n $. Consider the $ \mathfrak{m}_A $-adic filtration\[
        \varOmega^{\bullet,\mathrm{an}}_{X/S}\supseteq\mathfrak{m}_A\varOmega^{\bullet,\mathrm{an}}_{X/S}\supseteq\mathfrak{m}_A^2\varOmega^{\bullet,\mathrm{an}}_{X/S}\supseteq\cdots\supseteq\mathfrak{m}_A^N\varOmega^{\bullet,\mathrm{an}}_{X/S}=0
    \]of $ \varOmega^{\bullet,\mathrm{an}}_{X/S} $, which is biregular, and denote the associated graded object by\[
        \mathrm{gr}\, \varOmega^{\bullet,\mathrm{an}}_{X/S}=\bigoplus \mathrm{gr}_i\varOmega^{\bullet,\mathrm{an}}_{X/S},
    \]where $ \mathrm{gr}_i\varOmega^{\bullet,\mathrm{an}}_{X/S}:=\mathfrak{m}_A^i\varOmega^{\bullet,\mathrm{an}}_{X/S}\Big/\mathfrak{m}_A^{i+1}\varOmega^{\bullet,\mathrm{an}}_{X/S} $. Notice that\[
        \mathrm{gr}_0\varOmega^{\bullet,\mathrm{an}}_{X/S}=\varOmega^{\bullet,\mathrm{an}}_{X/S}\Big/\mathfrak{m}_A\varOmega^{\bullet,\mathrm{an}}_{X/S}=\varOmega^{\bullet,\mathrm{an}}_{X_0/\mathbb{C} }
    \]gives a acylic resolution of the constant sheaf $ \mathbb{C}  $, where $ X_0 $ is the smooth fiber of $ \mathrm{Spec}\, \mathbb{C} \hookrightarrow \mathrm{Spec}\, A $. Then\[
        \mathrm{gr}_{\mathfrak{m}_A}A\otimes_{\mathbb{C} }\mathrm{gr}_0\varOmega^{\bullet,\mathrm{an}}_{X/S}=\bigoplus_{n\geq 0}\mathfrak{m}^n_A/\mathfrak{m}^{n+1}_A\otimes_{A/\mathfrak{m}_A}\varOmega^{\bullet,\mathrm{an}}_{X/S}\Big/\mathfrak{m}_A \varOmega^{\bullet,\mathrm{an}}_{X/S}=\bigoplus_{n\geq 0}\mathfrak{m}_A^n\varOmega^{\bullet,\mathrm{an}}_{X/S}\Big/\mathfrak{m}_A^{n+1}\varOmega^{\bullet,\mathrm{an}}_{X/S}=\mathrm{gr}_{\mathfrak{m}_A}\varOmega^{\bullet,\mathrm{an}}_{X/S}
    \]resolves $ \mathrm{gr}_{\mathfrak{m}_A}A $, and hence $ \varOmega^{\bullet,\mathrm{an}}_{X/S} $ gives a acylic resolution of the constant sheaf $ A $. Since $ X $ is proper, by the GAGA principle we have\[
        \mathsf{R}^n\Gamma(X,\varOmega^{\bullet}_{X/S})\cong \mathsf{R}^n\Gamma(X^{\mathrm{an}},\varOmega^{\bullet,\mathrm{an}}_{X/S})=H^n(X^{\mathrm{an}},A)=H^n(X^{\mathrm{an}},\mathbb{C} )\otimes_{\mathbb{C} }A,
    \]and thus\begin{equation}\label{gio}
        \mathrm{length}_{A}\mathsf{R}^n\Gamma(X,\varOmega^{\bullet}_{X/S})=\mathrm{length}_A(A)\cdot\dim_{\mathbb{C} }H^n(X^{\mathrm{an}},\mathbb{C} ).
    \end{equation}
    \begin{lemma}\label{vvv}
        Let $ Y=\mathrm{Spec}\, A $ be a noetherian affine scheme, and let $ f:X\rightarrow Y $ be a proper morphism. For any complex $ \mathscr{F}^{\bullet}\in \mathsf{D}^{+}(\mathrm{Coh}(X)) $ that is flat over $ Y $, there exists a complex $ \mathscr{L}^{\bullet}\in\mathsf{D}^{+}(\mathrm{Coh}(Y)) $, whose terms $ \mathscr{L}^i $ has the form $ \mathcal{O}^{m_i}_Y $, and a functorial quasi-isomorphism\[
            f_*(\mathscr{F}^{\bullet}\otimes_{\mathcal{O}_X} f^*\mathscr{G})\longrightarrow\mathscr{L}^{\bullet}\otimes_{\mathcal{O}_Y}\mathscr{G}
        \]for any $ \mathscr{G}\in\mathrm{Qcoh}(Y) $. Moreover, for arbitrary base change diagram% https://quiver.taskfirst.cn/#q=WzAsNCxbMCwwLCJYJyJdLFsxLDAsIlgiXSxbMCwxLCJZJyJdLFsxLDEsIlkiXSxbMCwxLCJ1JyJdLFswLDIsImYnIl0sWzEsMywiZiJdLFsyLDMsInUiXV0=
        \[\begin{tikzcd}
            {X'} & X \\
            {Y'} & Y
            \arrow["{u'}", from=1-1, to=1-2]
            \arrow["{f'}", from=1-1, to=2-1]
            \arrow["f", from=1-2, to=2-2]
            \arrow["u", from=2-1, to=2-2]
        \end{tikzcd}\]there is a natural quasi-isomorphism\[
            f'_*(u'^*\mathscr{F}^{\bullet}\otimes_{\mathcal{O}_{X'}}f'^*\mathscr{G}')\longrightarrow u^*\mathscr{L}^{\bullet}\otimes_{\mathcal{O}_{Y'}}\mathscr{G}'
        \]for any $ \mathscr{G}'\in \mathrm{Qcoh}(Y') $.
    \end{lemma}
    
    \begin{proof}
        \cite[Prop 6.10.5]{EGA3_2}
    \end{proof}
    
    Since $ f $ is smooth, the sheaf $ \varOmega_{X/S} $ is locally free on $ X $. Applying (\ref{vvv}) to the pair $ (\mathscr{F}^{\bullet},\mathscr{G})=(\varOmega^{p}_{X/S},\mathcal{O}_Y) $ (Here we consider $ \varOmega^p_{X/S} $ as a complex concentrated in degree 0) and the base change morphism $ u:\mathrm{Spec}\, \mathbb{C} \rightarrow\mathrm{Spec}\, A $, we obtain\[
        \mathsf{R}f_*\varOmega^{p}_{X/S}\cong \mathscr{L}^{\bullet},\ \mathsf{R}f'_*\varOmega^{p}_{X_0/\mathrm{Spec}\, \mathbb{C} }\cong\mathsf{R}f_*'(u'^*\varOmega^{p}_{X/S})\cong u^*\mathscr{L}^{\bullet}.
    \]Therefore\[
        H^q(X,\varOmega^p_{X/S})=H^0(\mathrm{Spec}\, A,\mathsf{R}^q f_*\varOmega^p_{X/S})=\Gamma(\mathrm{Spec}\, A,H^q(\mathscr{L}^{\bullet}))
    \]is finite generated $ A $-module, and the natural isomorphism of base change\[
        H^q(X,\varOmega^p_{X/S})\otimes_A \mathbb{C} \cong\Gamma(\mathrm{Spec}\, \mathbb{C} ,H^q(u^*\mathscr{L}^{\bullet}))=H^q(X_0,\varOmega^p_{X_0}).
    \]
    
    \begin{lemma}
        Let $ (A,\mathfrak{m}_A) $ be the noetherian local ring with residue field $ k $. For any finite generated $ A $-module $ M $, we have the inequality\[
            \mathrm{length}_A(M)\leq\mathrm{length}_A(A)\cdot\dim_k(M\otimes_A k).
        \]Equality holds if and only if $ M $ is free $ A $-module.
    \end{lemma}

    \begin{proof}
        Let $ d=\dim_k(M\otimes_A k) $. By Nakayama's lemma, there exists a surjection $ f:A^d\twoheadrightarrow M $. Then\[
            \mathrm{length}_A(M)=d\cdot\mathrm{length}_A(A)+\mathrm{length}_A(\ker f)
        \]and $ \mathrm{length}_A(\ker f)=0 $ if and only if $ \ker f=0 $.
    \end{proof}

    By the above lemma, we have\begin{equation}\label{giu}
        \mathrm{length}_A H^q(X,\varOmega^p_{X/S})\leq\mathrm{length}_A(A)\cdot\dim H^q(X_0,\varOmega^p_{X_0})
    \end{equation}
    and equality holds if and only if $ H^q(X,\varOmega^p_{X/S}) $ is free $ A $-module. Consider Hodge to de Rham spectral sequence\[
        E^{p,q}_1=H^q(X,\varOmega^p_{X/S})\Longrightarrow \mathsf{R}^{p+q}\Gamma(\varOmega^{\bullet}_{X/S}).
    \]Since $ E^{p,q}_{\infty} $ is the subquotient of $ E^{p,q}_1 $, we deduce that\begin{equation}\label{giy}
        \mathrm{length}_A\mathsf{R}^n\Gamma(\varOmega^{\bullet}_{X/S})=\sum_{p+q=n}\mathrm{length}_A E^{p,q}_{\infty}\leq \sum_{p+q=n}\mathrm{length}_A H^q(X,\varOmega^p_{X/S}).
    \end{equation}
    Finally, combining (\ref{gio}), (\ref{giu}), and (\ref{giy}), we obtain that\begin{equation}
        \begin{aligned}\label{rrr}
            \mathrm{length}_A(A)\cdot\dim_{\mathbb{C} }H^n(X^{\mathrm{an}},\mathbb{C} )&=\mathrm{length}_A\mathsf{R}^n\Gamma(\varOmega^{\bullet}_{X/S})\\ &\leq\sum_{p+q=n}\mathrm{length}_A H^q(X,\varOmega^{p}_{X/S})\\ &\leq\mathrm{length}_A(A)\sum_{p+q=}\dim H^q(X_0,\varOmega^p_{X_0}).
        \end{aligned}
    \end{equation}
    By Hodge decomposition in analytic setting ($ X_0 $ is projective so $ X_0^{\mathrm{an}} $ is K\"{a}hler) and the GAGA principle, we have\[
        H^n(X^{\mathrm{an}},\mathbb{C} )=H^n(X_0^{\mathrm{an}},\mathbb{C} )\cong\bigoplus_{p+q=n}H^q(X_0^{\mathrm{an}},\varOmega^{p,\mathrm{an}}_{X_0})\cong\bigoplus_{p+q=n}H^q(X_0,\varOmega^{p}_{X_0}),
    \]and hence equality holds throughout in (\ref{rrr}). It follows that each\[
        H^q(X,\varOmega^p_{X/S})
    \]is a free $ A $-module.
\end{proof}

\begin{theorem}
    Let $ X $ be a smooth projective variety over $ \mathbb{C}  $. If $ \omega_X\cong\mathcal{O}_X $, then the flat deformations of $ X $ (cf. \ref{B8}) are unobstructed.
\end{theorem}

\begin{proof}
    By (\ref{qqq}), it suffices to show that\[
        \mathrm{Def}(X_{n+1}/A_{n+1},B_{n+1})\longrightarrow\mathrm{Def}(X_n/A_n,B_n)
    \]is surjective for all positive integer $ n $ and $ X_n=F(\alpha_n)(X_{n+1}) $. By (\ref{iop}), the problem becomes to show that\begin{equation}\label{afg}
        H^1(X_{n+1},\mathcal{T}_{X_{n+1}/A_{n+1}})\longrightarrow H^1(X_n,\mathcal{T}_{X_n/A_n})
    \end{equation}is surjective.

    Let $ f:X_n\rightarrow \mathrm{Spec}\, A_n $ be the deformation of $ X $ over $ A_n $, which is projective. By (\ref{kkk}), it follows that\[
        H^q(X_n,\varOmega^p_{X_n/A_n})
    \]is a free $ A_n $-module of rank $ h^{p,q}_n $. By cohomology and base change, we have the isomorphisms\[
        H^q(X_n,\varOmega^p_{X_n/A_n})\otimes_{A_n}k\stackrel{\sim}{\longrightarrow}H^q(X_0,\varOmega^p_X),
    \]and thus $ h^{p,q}_n=h^{p,q}(X) $ is independent on $ n $. In particular,\[
        H^0(X,\omega_X)=k\Longrightarrow H^0(X_n,\omega_{X_n/A_n})=A_n
    \]Take a non-trivial section $ s:\mathcal{O}_{X_n}\rightarrow\omega_{X_n/A_n} $, then $ s\otimes_{X_n}X $ is an isomorphism on $ X $. Since $ \omega_{X_n/A_n} $ is flat over $ \mathrm{Spec}\, A_n $, by (\ref{ghu}), $ s $ is an isomorphism. From the perfect pairing\[
        \varOmega^1_{X_n/A_n}\otimes\varOmega^{d-1}_{X_n/A_n}\longrightarrow\omega_{X_n/A_n}\cong\mathcal{O}_{X_n},
    \]we have $ \mathcal{T}_{X_n/A_n}\cong\varOmega^{d-1}_{X_n/A_n} $. Now we may rewrite (\ref{afg}) to\[
        A^{h^{d-1,1}}_{n+1}\longrightarrow A^{h^{d-1,1}}_n,
    \]which is a surjection as desired.
\end{proof}

\section{Threefold Calabi-Yau varieties with ordinary double points}

We summarize classical results concerning the deformation theory of threefold Calabi-Yau varieties with ordinary double points (ODP).
In particular, we state Friedman's criterion for smoothability and the later generalization by Kawamata, Ran and Tian, which shows that deformations are unobstructed once a K\"ahler resolution exists.

\subsection{Setup}
Let $X_0$ be a three-dimensional projective variety with $\omega_{X_0}\cong\mathcal{O}_{X_0}$ and only ordinary double points as singularities.
A \emph{small resolution} is a proper birational morphism $f\colon X\to X_0$ such that $X$ is smooth and the exceptional set is a union of rational curves $C_1,\dots,C_r\cong\mathbb{P}^1$ with normal bundle $\mathcal{O}_{C_i}(-1)\oplus\mathcal{O}_{C_i}(-1)$.
We denote by $[C_i]$ the homology class in $H_2(X,\mathbb{C})$ or its dual in $H^4(X,\mathbb{C})$.
The variety $X_0$ is called \emph{smoothable} if it occurs as the central fiber of a flat family $X\to\Delta$ with smooth total space and smooth nearby fibers $X_t$ (here $\Delta$ is a small disc).

\subsection{Friedman's theorem}
The first precise criterion for smoothability was given by Friedman~\cite{Friedman1986}.

\begin{theorem}[Friedman, 1986]
Let $X_0$ be as above and let $f\colon X\to X_0$ be a small resolution.
Then $X_0$ is smoothable (and in fact the small resolution can be deformed to an isomorphism) if the following two conditions hold:
\begin{enumerate}
    \item The classes $[C_1],\dots,[C_r]$ generate the cohomology group $H^2(X,\Omega_X^2)$.
    \item There exists a relation $\sum_{i=1}^r a_i[C_i]=0$ in $H^4(X,\mathbb{C})$ with all coefficients $a_i\neq 0$.
\end{enumerate}
\end{theorem}

\subsection{Unobstructedness and the theorem of Kawamata-Ran-Tian}
Subsequently it was shown independently by Y.\ Kawamata, Z.\ Ran and G.\ Tian~\cite{Kawamata,Ran,Tian} that the first generating condition is superfluous provided that a K\"ahler resolution exists.
The crucial point is a generalization of the Bogomolov-Tian-Todorov unobstructedness theorem to singular Calabi-Yau varieties.

\begin{theorem}[Kawamata, Ran, Tian]
Let $X_0$ be a threefold with $\omega_{X_0}\cong\mathcal{O}_{X_0}$ and only ordinary double points.
Assume there exists a resolution $p\colon X'\to X_0$ such that $X'$ is a compact K\"ahler manifold.
Then the versal deformation space $\operatorname{Def}(X_0)$ is smooth, i.e.\ deformations of $X_0$ are unobstructed.
In particular, the condition that the exceptional curves generate $H^2(\Omega_X^2)$ in Friedman's theorem can be dropped.
\end{theorem}

\subsection{Sketch of proofs}

\subsubsection{Friedman's smoothing criterion}

\begin{enumerate}
    \item \textbf{Deformations of a small resolution.}
    Let $f\colon X\to X_0$ be a small resolution with exceptional curves $C_i$.
    A deformation of $X$ simultaneously blows down the $C_i$ if the associated infinitesimal deformation maps the normal bundle sequence so that the curves are obstructed from deforming.
    More precisely, the local-to-global exact sequence for $\operatorname{Ext}^1$ and $\operatorname{Ext}^2$ relates the first-order deformations of $X$ and of the formal completion along each $C_i$.
    Condition~\emph{(a)} (generating $H^2(\Omega_X^2)$) guarantees that the natural map $\operatorname{Ext}^1(\Omega_X^1,\mathcal{O}_X)\to \bigoplus_i \operatorname{Ext}^1(\mathcal{O}_{C_i},\mathcal{O}_{C_i}(-1)^{\oplus 2})$ is surjective, so one can lift local smoothings of the double points to a global deformation.
    
    \item \textbf{Simultaneous contraction and the relation among $[C_i]$.}
    In order to contract all $C_i$ simultaneously in a family, the exceptional set must become trivial in homology.
    The obstruction to contracting a curve in a deformation is its fundamental class; a family of threefolds containing $X$ as a fiber and for which the curves disappear corresponds to a first-order deformation $\nu\in H^1(X,\Theta_X)$ for which the restriction to each $C_i$ vanishes.
    This translates precisely to the existence of a relation $\sum a_i[C_i]=0$ in $H^4(X,\mathbb{C})$ with every $a_i\neq 0$.
    The coefficients $a_i$ determine the ``speed'' at which each singularity is smoothed.
    
    \item \textbf{Conclusion.}
    Once one has such a first-order deformation, the unobstructedness (proved in~\ref{wri}) ensures that it can be lifted to an actual smoothing of $X_0$.
    The work of Kawamata-Ran-Tian shows that the extra generating condition is automatically satisfied when a K\"ahler resolution exists, completing the picture.
\end{enumerate}

\subsubsection{Unobstructedness for ODP Calabi-Yau threefolds}\label{wri}

\begin{enumerate}
    \item \textbf{Deformation functor and tangent-obstruction theory.}
    Since $X_0$ is a local complete intersection (ODP), its deformation theory is controlled by the cotangent complex $\mathbb{L}_{X_0}$.
    The tangent space is $\operatorname{Ext}^1(\Omega_{X_0}^1,\mathcal{O}_{X_0})$ and obstructions lie in $\operatorname{Ext}^2(\Omega_{X_0}^1,\mathcal{O}_{X_0})$.
    Smoothness of $\operatorname{Def}(X_0)$ follows if $\operatorname{Ext}^2(\Omega_{X_0}^1,\mathcal{O}_{X_0})=0$.
    
    \item \textbf{Comparison with a K\"ahler resolution.}
    Let $p\colon X'\to X_0$ be a K\"ahler resolution (not necessarily small).
    The sheaf of differentials $\Omega_{X_0}^1$ sits in an exact triangle relating it to $p_*\Omega_{X'}^1$ and a sheaf supported on the singular points.
    Using local calculations one shows the connectedness of the deformation functors, namely $\operatorname{Def}(X_0)\cong \operatorname{Def}(X')$ near the central fiber.
    Hence it is enough to prove $\operatorname{Ext}^2(\Omega_{X'}^1,\mathcal{O}_{X'}) =0$ for the smooth K\"ahler threefold $X'$.

    \item \textbf{Hodge theory and the Bogomolov--Tian--Todorov argument.}
    On the smooth compact K\"ahler manifold $X'$ we have the Hodge decomposition and Serre duality.
    Because $p$ contracts only rational curves, one still has $\omega_{X'}\cong\mathcal{O}_{X'}$.
    The classical unobstructedness proof for Calabi--Yau manifolds (due to Bogomolov, Tian, Todorov) uses the contraction map
    \[
        H^2(X',\Theta_{X'})\times H^1(X',\Theta_{X'})\to H^3(X',\Theta_{X'})
    \]
    given by the cup product with the trace map $\Theta_{X'}\otimes\Omega_{X'}^p\to \Omega_{X'}^{p-1}$.
    Under the Calabi--Yau condition $\omega_{X'}\cong\mathcal{O}_{X'}$, Serre duality identifies the obstruction space $\operatorname{Ext}^2(\Omega_{X'}^1,\mathcal{O}_{X'})\cong H^2(X',\Theta_{X'})$ with the dual of $H^1(X',\Omega_{X'}^2)$.
    By a clever use of the local Torelli theorem for the period map, one shows that the multiplication map vanishes, hence obstructions vanish.
    This argument adapts to the present situation and yields $\operatorname{Ext}^2(\Omega_{X_0}^1,\mathcal{O}_{X_0})=0$.
\end{enumerate}


% Now let us comparing the functors $ F $ and $ F_1 $ mentioned at the beginning. Recall that given a scheme $ X_0/k $, $ F $ is the fucntor of deformations of $ X_0 $ over $ A $, that is, pairs $ (X,i) $ where $ X $ is flat over $ A $, and $ i:X_0\hookrightarrow X $ is a morphism such that $ i\otimes k:X_0 \rightarrow X\otimes k $ is an isomorphism, while $ F_1 $ is the functor of flat families $ X/A $ such that there exists an isomorphism $ X\otimes k\cong X_0 $. There is a natural `forgetful' fucntor from $ F $ to $ F_1 $.

% The following result is proved using the same kind of arguments as in the earlier part of this subsection.

% \begin{theorem}
%     Suppose the hypothesis of (\ref{yui}) satisfied. Then:
%     \begin{enumerate}[(1)]
%         \item 
%     \end{enumerate}
% \end{theorem}

% \begin{lemma}\label{ghu}
%     Let $ A $ be a ring, $ J $ be a nilpotent ideal in $ A $, and $ u:M\rightarrow N $ be a homomorphism of $ A $-modules, with $ N $ falt over $ A $. If $ \bar{u}:M/JM\rightarrow N/JN $ is an isomorphism, then $ u $ is an isomorphism.
% \end{lemma}

% \begin{proof}
%     Let $ K=\mathrm{coker}\, u $. Tensor the exact sequence\[
%         M\rightarrow N\rightarrow K\rightarrow 0
%     \]with $ A/J $, we have $ K/JK=0 $. Since $ J $ is nilpotent, we have $ K=0 $. If $ K'=\ker u $, then we get an exact sequence\[
%         0\rightarrow K'/JK'\rightarrow M/JM\rightarrow N/JN\rightarrow 0
%     \]by the flatness of $ N $. Hence, $ K'=0 $.
% \end{proof}

% Fix a scheme $ X $ over $ k $.

% \begin{definition}
%     An \textbf{(infinitesimal) deformation} of $ X $ over $ A\in\mathsf{C} $ is a flat scheme $ Y $ over Spec$ A $ with an immersion $ i:X\hookrightarrow Y $ inducing an isomorphism $ X\stackrel{\sim}{\longrightarrow}Y\times_A k $
%     % https://quiver.taskfirst.cn/#q=WzAsNSxbMCwwLCJYIl0sWzAsMSwiXFxtYXRocm17U3BlY31cXCwgayJdLFsxLDAsIlkiXSxbMSwxLCJcXG1hdGhybXtTcGVjfVxcLCBBIl0sWzIsMCwiXFxidWxsZXQiXSxbMCwxXSxbMCwyLCJpIl0sWzEsMywiIiwyLHsibGFiZWxfcG9zaXRpb24iOjB9XSxbMiwzXV0=
%     \[\begin{tikzcd}
%         X & Y \\
%         {\mathrm{Spec}\, k} & {\mathrm{Spec}\, A.}
%         \arrow["i", from=1-1, to=1-2]
%         \arrow[from=1-1, to=2-1]
%         \arrow[from=1-2, to=2-2]
%         \arrow[from=2-1, to=2-2]
%     \end{tikzcd}\]
%     If $ Y' $ is another deformation to $ A $, then $ Y $ and $ Y' $ are \textbf{isomorphic} if there exists a morphism $ f:Y\rightarrow Y' $ over $ A $ which induces the identity on the closed fiber $ X $. By~\ref{ghu} $ f $ must be an isomorphism of schemes.
% \end{definition}

% Define the deformation functor $ \mathcal{D}=\mathrm{Def}_X $ by\[
%     \mathcal{D}(A)=\{\text{isomorphism classes of deformations of}\ X/k\ \text{to}\ A\}.
% \]
% Unfortunately, $ \mathcal{D} $ is not pro-representable. But in some finiteness restrictions on $ X $, $ \mathcal{D} $ will have a hull.

% Suppose that $ (A',\eta')\rightarrow (A,\eta) $ and $ (A'',\eta'')\rightarrow (A,\eta) $ are morphisms of couples, where $ A''\rightarrow A $ is a surjection. Let $ Y',Y,Y'' $ denote deformations in the class of $ \eta',\eta,\eta'' $, respectively. The sheaf $ \mathcal{O}_{Y'}\times_{\mathcal{O}_Y}\mathcal{O}_{Y''} $ of $ A'\times_A A'' $-algebras defined a scheme $ Z $ on $ X $, which is flat over $ B:=A'\times_A A'' $. Moreover, $ Z\in \mathcal{D}(B) $ maps to $ (\eta',\eta'')\in \mathcal{D}(A')\times_{\mathcal{D}(A)}\mathcal{D}(A'') $. Hence $ \mathcal{D} $ satisfies the condition $ H_1 $.

% Suppose that $ W $ is another deformation over $ B $ with morphism of deformation

\newpage

\bibliographystyle{unsrt}
\bibliography{reference}

\end{document}